% !TeX root = ../../Book4.tex
% 纲要标题：## 附录 C　连续上同调与几何应用
% 纲要位置：工作记录/计划纲要.md，第 3898 行。
% 纲要细目（写作范围）：
% > 以离散连续系数为固定范围，证明连续余链与导出函子的识别并计算低次例子；层上同调与有限 Čech 覆盖给出比较证明，拓扑比较、仿射消失、平展层比较及局部域 Brauer 不变量给出证明概要。


\chapter{连续上同调与几何应用}
\label{b4:app:C}

\begin{chapterabstract}
本附录分别研究投射有限群的连续上同调与有限覆盖下的层上同调，证明连续余链的有限商计算和 Čech 比较，并给出 Galois 与射影直线上的例子。最后从有限连续作用的纤维函子重建投射有限群。
\end{chapterabstract}

\begin{learninggoals}
\begin{itemize}
\item 区分抽象群模、离散连续模与带非离散拓扑的系数，说明有限商计算中指标和映射的方向。
\item 用连续交叉同态计算一次上同调，用膨胀映射判断有限阶段的二次上同调类是否消失。
\item 证明连续 Hilbert 第 90 定理，保留 Kummer 系数的自然 Galois 作用，并计算 \(\mathbb F_q(\!(t)\!)\) 的素于特征的幂商。
\item 写出层的截面、茎与内射消解，建立 Čech 双复形，核验零调覆盖的比较条件。
\item 通过具体覆盖区分固定覆盖的 Čech 上同调、层上同调及域上的平展上同调，逐项检查比较定理所需的空间与系数条件。
\item 用正则有限商上的自然性恢复纤维函子的全部自同构。
\end{itemize}
\end{learninggoals}

有限群上同调用函数 \(G^n\rightarrow A\) 写余链。若 \(G\) 换成投射有限群，而 \(A\) 保持离散，同一个公式还须加上连续性；这迫使一条余链在足够小的开闭陪集上取常值，于是计算可以降到某个有限商。几何中也有相似的局部到整体问题：若分别在几个开集上给出截面，要把它们粘成整体截面，就须检查交叠处的值是否一致。两类问题使用的空间不同，却都从「局部可写的数据何时相容」出发，并把失败记录为上同调类。本附录分别在连续群余链与有限覆盖的 \v{C}ech 余链\footnote{切赫（Eduard Čech，1893-1960），捷克数学家，研究拓扑学，发展以覆盖定义的上同调方法。}中完成这些计算。

\ProseParagraphBreak
本附录使用\chapref{b4:ch:11}的群上同调、\cref{b4:prop:filtered-colimits-exact}的滤过余极限正合性，以及\cref{b4:cor:double-complex-comparison}的双复形比较，还用到有限群与有限 Galois 扩张的基础知识。需要的连续系数条件和层的基本构造在此重述。射影直线例子只涉及两张仿射坐标图与 Laurent 多项式，其层上同调解释采用一项仿射消失定理。本附录不系统讲授一般概形、局部类域论或平展景理论。

\ProseParagraphBreak
本附录的两部分各有自己的预备知识。\secref{b4:sec:C01}承接\chapref{b4:ch:11}的群上同调，\secref{b4:sec:C02}发展连续余圈和 Galois 算例；其中局部域的计算另需基本赋值与有限域知识。

\ProseParagraphBreak
\secref{b4:sec:C03}直接从\chapref{b4:ch:08}的右导出函子出发，与连续群的两节无依赖关系。当有限开覆盖的每个非空有限交集上，所论交换群层的正次层上同调全为零时，该节证明覆盖的 Čech 上同调计算层上同调；证明还用到\secref{b4:sec:0604}的双复形工具。\secref{b4:sec:C04}转向拓扑比较和平展层，需要相应拓扑或概形基础，其中的比较结果按正文标明的证明状态使用。

% !TeX root = ../../Book4.tex
% 纲要标题：### C.1 投射有限群与连续离散模
% 纲要位置：工作记录/计划纲要.md，第 3902 行。
% 纲要细目（写作范围）：
% #### C.1.1 有限商与开正规子群
% #### C.1.2 离散连续模与连续余链
% #### C.1.3 有限商计算与导出函子


\section{投射有限群与连续离散模}
\label{b4:sec:C01}

有限商决定投射有限群的拓扑，但一个连续余链所经过的有限商会随余链变化。连续上同调因此需要把所有有限阶段及其过渡映射一起考虑。本节从\cref{b4:def:continuous-group-cochains}的系数范畴出发，将连续上同调识别为不变量函子的右导出，并说明\cref{b4:prop:continuous-finite-quotients}的有限商计算。

\subsection{有限商与开正规子群}
\label{b4:subsec:C01:01}

投射有限群 \(G\) 是有限离散群的投射极限，取有限群直积的子空间拓扑。群运算逐坐标进行，故连续；相容条件定义闭子集，因此 \(G\) 紧且 Hausdorff。有限多个坐标投影核的交是开正规子群，并组成单位元处的邻域基。这里所需的紧性是有限离散空间直积的紧性；同一拓扑事实已在\subsecref{b4:subsec:1104:03}说明。

\ProseParagraphBreak
开子群 \(H\le G\) 的各陪集覆盖紧空间 \(G\)，因而只有有限多个，故 \([G:H]<\infty\)。它也是闭集，因为补集是其他开陪集之并。取它的有限多个共轭子群之交，便得到包含在 \(H\) 中的开正规子群。反过来，若闭子群有有限指数，则其每个陪集的补集是有限个闭陪集之并，故它是开子群。这些判断中不能把“闭”省掉后随意把抽象有限指数子群当作开子群。

\begin{proposition}\label{b4:prop:C-finite-module-quotient}
设 \(G\) 为投射有限群，\(M\) 为有限交换群。\(G\) 对离散 \(M\) 的作用连续，当且仅当存在开正规子群 \(U\triangleleft G\)，使 \(U\) 在 \(M\) 上平凡作用；此时作用通过有限群 \(G/U\)。此外，从 \(G\) 到有限离散群的群同态连续，当且仅当其核为开子群。
\end{proposition}
\begin{proof}
连续作用使每个 \(m\in M\) 的稳定子群开；将有限多个稳定子群相交，得到作用同态的核，这是开正规子群。反过来，经有限离散商的作用显然连续。对于群同态，各个非空纤维都是核的陪集，故所有纤维开与核开等价，也就与到离散群的连续性等价。
\end{proof}

两类基本例子是
\[
 \widehat{\vphantom{Z}\mathbb Z}=\varprojlim_n\mathbb Z/n\mathbb Z,
 \qquad
 \mathbb Z_p=\varprojlim_r\mathbb Z/p^r\mathbb Z.
\]
第一式按正整数的整除关系取极限，第二式固定素数 \(p\)。整数的像在两者中都稠密，\(1\) 是一个拓扑生成元；但 \(\widehat{\vphantom{Z}\mathbb Z}\) 具有任意阶的有限循环商，\(\mathbb Z_p\) 的有限连续商只能是循环 \(p\)-群。后一限制将直接反映在上同调计算中。
\symidx{\(\widehat{\vphantom{Z}\mathbb Z}\)：整数群的投射有限完备化}

\subsection{离散连续模与连续余链}
\label{b4:subsec:C01:02}

设 \(G\) 为投射有限群。一个离散连续 \(G\)-模是带 \(G\) 作用的交换群 \(M\)，其中每个元素的稳定子群都是开的；这正是\cref{b4:def:continuous-group-cochains}中作用映射连续的等价条件。每个开稳定子群又包含一个开正规子群，所以这里的条件是：对每个 \(m\in M\)，存在固定 \(m\) 的开正规子群 \(U_m\)。这个子群可以随 \(m\) 改变，并不要求存在一个开正规子群固定整个 \(M\)。\cref{b4:prop:C-no-common-action-quotient}将给出没有共同有限作用商的例子。对象不必有限。例如，若 \(K^s\) 是域 \(K\) 的可分闭包，则每个代数元被某个有限扩张的 Galois 子群固定，所以 \((K^s,+)\) 和 \((K^s)^\times\) 都是离散连续 \(G_K\)-模，其中 \(G_K=\operatorname{Gal}(K^s/K)\) 取 Krull 拓扑。

\ProseParagraphBreak
记这些模及 \(G\)-同态所成的范畴为 \(\mathcal D_G\)。核、像和商都仍为离散连续模：子模中的元素保留原稳定子群，商中的元素则由任一代表的开稳定子群固定。因此 \(\mathcal D_G\) 是交换范畴，嵌入全部抽象 \(\mathbb ZG\)-模的函子正合。不过，两范畴的内射对象不能不经检查便混为一谈。

\ProseParagraphBreak
连续余链为 \(C^n_{\mathrm{cts}}(G,M)=\operatorname{Map}_{\mathrm{cts}}(G^n,M)\)，次数零取 \(M\)。为方便独立查阅，非齐次微分写为
\begin{align*}
 (\delta f)(g_1,\ldots,g_{n+1})
 &=g_1f(g_2,\ldots,g_{n+1})\\
 &\quad+\sum_{i=1}^n(-1)^i
 f(g_1,\ldots,g_ig_{i+1},\ldots,g_{n+1})\\
 &\quad+(-1)^{n+1}f(g_1,\ldots,g_n).
\end{align*}
零次公式为 \((\delta m)(g)=gm-m\)。这是\eqref{b4:eq:group-cochain-differential}的同一公式，微分平方为零的代数验证继续适用；连续性则来自乘法、作用与有限求和的连续性。

\begin{proposition}\label{b4:prop:C-cochains-exact}
设 \(G\) 为投射有限群，\(\mathcal D_G\) 为离散连续 \(G\)-模及 \(G\)-同态组成的交换范畴。对每个 \(n\ge0\)，函子 \(C^n_{\mathrm{cts}}(G,-):\mathcal D_G\rightarrow\mathbf{Ab}\) 正合。因此连续群上同调对离散连续模的短正合列给出长正合列。
\end{proposition}
\symidx{\(\mathcal D_G\)：离散连续群作用模的范畴}
\begin{proof}
核和单射逐点计算。对满射 \(M\rightarrow N\)，一个连续余链 \(f:G^n\rightarrow N\) 从紧空间映到离散空间，像只有有限多个值。为每个值选一个 \(M\) 中的原像，在它的开闭纤维上取对应常值，得到连续提升。这只是在提升一张连续函数，不要求它等变，也不要求它满足余圈方程；所选原像也不必固定于原先同一个开子群。次数零就是原满射。于是短正合列给出逐次短正合的余链复形，长正合列来自\cref{b4:thm:cohomology-long-exact}的连接态射构造。
\end{proof}

余链可以提升，并不意味着余圈可以提升为余圈。若 \(0\to J\to M\to N\to0\) 中的余圈 \(f\) 提升为余链 \(\widetilde f\)，则 \(\delta\widetilde f\) 在 \(N\) 中的像为 \(\delta f=0\)，因而取值于 \(J\)；它的上同调类就是连接同态记录的提升障碍。次数零的余圈是不变元素，\cref{b4:prop:C-invariant-lift-obstruction}将具体展示一个元素有提升，却没有不变提升的情形。

\subsection{有限商计算与导出函子}
\label{b4:subsec:C01:03}

\cref{b4:prop:continuous-finite-quotients}已经证明，对任意离散连续 \(G\)-模 \(M\)，
\begin{equation}\label{b4:eq:C-finite-quotient-cohomology}
 H^n_{\mathrm{cts}}(G,M)
 \cong\varinjlim_{U\triangleleft_{\mathrm{open}}G}H^n(G/U,M^U).
\end{equation}
指标按开正规子群的反向包含排列：从 \(U\) 到更小的 \(V\) 使用膨胀及系数包含 \(M^U\subseteq M^V\)。证明中的两个有限性步骤是：连续余链像有限，并且它在某个共同开正规子群的陪集上常值。再缩小这个子群使其固定有限像，余链便通过 \((G/U)^n\) 且取值于 \(M^U\)。\cref{b4:prop:filtered-colimits-cohomology}保证取余极限与取上同调相容。

\ProseParagraphBreak
若 \(M\) 有限，先固定一个作用核 \(U_0\)，以后仅取 \(U\subseteq U_0\) 即可，这些指标共尾，且 \(M^U=M\)。但不能只计算 \(H^n(G/U_0,M)\) 后停止：更小子群带来的膨胀映射可能使原有上同调类消失。

\begin{theorem}\label{b4:thm:C-continuous-derived}
设 \(G\) 为投射有限群，\(\mathcal D_G\) 为离散连续 \(G\)-模及 \(G\)-同态组成的交换范畴。则 \(\mathcal D_G\) 有足够多内射对象，并且对每个 \(n\ge0\)，有
\[
 H^n_{\mathrm{cts}}(G,M)\cong R^n(-)^G(M)
 \qquad(M\in\mathcal D_G).
\]
右边的导出函子在离散连续模范畴内计算。
\end{theorem}
\begin{proof}
先从抽象群模范畴取得内射对象，再使其成为离散连续模。对任意抽象 \(\mathbb ZG\)-模 \(B\)，令 \(D(B)\) 为所有被某个开子群固定的元素所成子模。有限交保证加法封闭，开子群的共轭仍开保证 \(G\) 稳定。每个从离散连续模 \(M\) 到 \(B\) 的同态都落入 \(D(B)\)，给出自然双射
\[
 \operatorname{Hom}_{\mathbb ZG}(M,B)
 \cong\operatorname{Hom}_{\mathcal D_G}(M,D(B)).
\]
所以离散化函子 \(D\) 是正合嵌入 \(\mathcal D_G\rightarrow\mathbb ZG\text{-}\mathbf{Mod}\) 的右伴随。抽象模范畴有足够多内射对象；把 \(M\) 嵌入一个抽象内射模 \(J\)，其像落入 \(D(J)\)。\cref{b4:prop:right-adjoint-injectives}说明 \(D(J)\) 在 \(\mathcal D_G\) 中内射，于是得到足够多内射对象。

再检查离散连续内射对象在每个有限商上的高次上同调。令 \(I\) 是 \(\mathcal D_G\) 中的内射对象。对开正规子群 \(U\)，不变量 \(I^U\) 是 \(G/U\)-模。此处使用另一个伴随：沿商映射的膨胀函子与取 \(U\)-不变量满足
\[
 \operatorname{Hom}_{\mathcal D_G}
 (\operatorname{Inf}_{G/U}^G N,I)
 \cong\operatorname{Hom}_{\mathbb Z[G/U]}(N,I^U).
\]
因为左边的同态必须把 \(N\) 映入 \(I^U\)，这个双射由同一张底群映射给出。膨胀不改变底群而正合，故\cref{b4:prop:right-adjoint-injectives}使 \(I^U\) 成为内射 \(G/U\)-模。有限群上同调是抽象群模不变量函子的右导出，因而 \(H^n(G/U,I^U)=0\) 对 \(n>0\) 成立；代入\eqref{b4:eq:C-finite-quotient-cohomology}，得到 \(H^n_{\mathrm{cts}}(G,I)=0\)。第一个伴随建立足够多内射对象，这个伴随则把高次消失化为各有限商上的消失。

由\cref{b4:prop:C-cochains-exact}，连续上同调构成上同调 \(\delta\)-函子，零次为 \(M^G\)。刚证的内射嵌入使其正次可消去，故\cref{b4:thm:derived-universal}将它自然识别为不变量函子的右导出函子。
\end{proof}

这给出了\subsecref{b4:subsec:1104:03}中使用的导出识别。上述范畴与有限商证明可对照 Weibel\cite[Definition 6.11.8; Lemma 6.11.10; Theorem 6.11.13, pp. 210--213]{Weibel1994HomologicalAlgebra}。如果系数带有 \(p\)-进拓扑，例如把 \(\mathbb Z_p\) 本身作为拓扑模，就已离开 \(\mathcal D_G\)；相关连续上同调还要处理逆极限与系数拓扑，本附录不作无说明的替换。

\begin{proposition}\label{b4:prop:C-no-common-action-quotient}
设 \(p\) 为素数，\(G=\mathbb Z_p\)。对 \(r\ge1\)，令 \(M_r=\mathbb F_p[\mathbb Z/p^r\mathbb Z]\)，使 \(G\) 经模 \(p^r\) 商在正则基上作左平移。则离散交换群
\[
 M=\bigoplus_{r\ge1}M_r
\]
是离散连续 \(G\)-模，但其作用不通过任何有限商。
\end{proposition}
\begin{proof}
每个 \(m\in M\) 只有有限多个非零分量。若这些分量的指标都不超过 \(r_0\)，则开正规子群 \(p^{r_0}\mathbb Z_p\) 在这些分量上平凡作用，故固定 \(m\)。零元素被整个 \(G\) 固定，所以每个元素的稳定子群都开，作用连续。

在 \(M_r\) 上，非零的模 \(p^r\) 剩余类会移动正则基中零剩余类所对应的基向量，故作用核恰为 \(p^r\mathbb Z_p\)。于是整个 \(M\) 上的作用核为
\[
 \bigcap_{r\ge1}p^r\mathbb Z_p=\{0\}.
\]
该等式来自逆极限中的元素由全部模 \(p^r\) 坐标决定。因 \(\mathbb Z_p\) 无限，零子群不可能具有有限指数，也不是开子群。因此不存在一个固定整个 \(M\) 的开正规子群，忠实作用也不可能通过有限群分解。
\end{proof}

\begin{proposition}\label{b4:prop:C-invariant-lift-obstruction}
设 \(p\) 为素数，\(G=\mathbb Z_p\) 经商群 \(C_p\) 作用于 \(M=\mathbb F_p[C_p]\)。令 \(\epsilon:M\to\mathbb F_p\) 为系数和，\(J=\ker\epsilon\)，并给 \(\mathbb F_p\) 平凡作用。则离散连续模的短正合列
\[
 0\longrightarrow J\longrightarrow M
 \xrightarrow{\epsilon}\mathbb F_p\longrightarrow0
\]
在每个连续余链次数仍为短正合列，但 \(M^G\to\mathbb F_p\) 是零映射。连接同态将 \(1\) 送到非零类
\[
 [g\mapsto ge-e]\in H^1_{\mathrm{cts}}(G,J),
\]
其中 \(e\) 是群代数中单位元素所对应的基向量。
\end{proposition}
\begin{proof}
正则平移作用使不变向量的全部系数相等，所以
\[
 M^G=\mathbb F_p\cdot\sum_{h\in C_p}h.
\]
这个基向量和的系数和为 \(p=0\)，故不变量映射为零。连续余链上的短正合性由\cref{b4:prop:C-cochains-exact}给出。

用零次余链 \(e\) 提升 \(1\)，其微分是 \(g\mapsto ge-e\)。因 \(\epsilon(ge-e)=0\)，它取值于 \(J\)；又通过有限商 \(C_p\)，故连续。这正是连接同态的代表。若存在 \(j\in J\) 使 \(ge-e=gj-j\) 对所有 \(g\) 成立，则 \(e-j\) 是系数和为 \(1\) 的不变向量，与刚才的零映射结论矛盾。所以该类非零。提升余链后的不变性误差在此不能通过修改提升消除。
\end{proof}

% !TeX root = ../../Book4.tex
% 纲要标题：### C.2 连续群上同调与局部域
% 纲要位置：工作记录/计划纲要.md，第 3613 行。
% 纲要细目（写作范围）：
% #### C.2.1 一次上同调与连续交叉同态
% #### C.2.2 投射循环群的计算
% #### C.2.3 Galois 上同调与局部域算例


\section{连续群上同调与局部域}
\label{b4:sec:C02}

连续上同调在低次已有具体含义：零次是固定元素，一次衡量连续交叉同态偏离主交叉同态的程度。先在投射循环群上计算，再把同样的公式用于 Galois 群，可以看清哪些结论只靠有限商，哪些结论还需要域的算术结构。

\subsection{一次上同调与连续交叉同态}
\label{b4:subsec:C02:01}

\begin{proposition}\label{b4:prop:C-continuous-h1}
设 \(G\) 为投射有限群，\(M\) 为离散连续 \(G\)-模。则 \(H^0_{\mathrm{cts}}(G,M)=M^G\)，且
\[
 H^1_{\mathrm{cts}}(G,M)
 =\frac{\bigl\{\,f:G\rightarrow M\text{ 连续}:f(gh)=f(g)+gf(h)\,\bigr\}}
 {\{\,g\mapsto gm-m:m\in M\,\}}.
\]
特别地，若作用平凡，则 \(H^1_{\mathrm{cts}}(G,M)=\operatorname{Hom}_{\mathrm{cts}}(G,M)\)，其中同态的目标使用加法。
\end{proposition}
\begin{proof}
零次余圈条件是 \(gm-m=0\)。一次余圈条件由微分给出
\(gf(h)-f(gh)+f(g)=0\)，零次余边正是 \(g\mapsto gm-m\)。作用平凡时余圈条件成为加法同态条件，全部一次余边为零。
\end{proof}

分子中的映射称为连续交叉同态；本卷\secref{b4:sec:1102}已给出其离散版本。连续条件在这里不能省略。例如，对平凡作用的离散 \(\mathbb Z\)，紧性使连续同态的像为有限子群，而 \(\mathbb Z\) 没有非零有限子群，所以任意投射有限群都满足 \(H^1_{\mathrm{cts}}(G,\mathbb Z)=0\)。

\ProseParagraphBreak
若系数用乘法写作 \(A\)，公式相应变为 \(f(gh)=f(g)\,g\bigl(f(h)\bigr)\)，主交叉同态为 \(g\mapsto g(a)/a\)。这只是记法改变；下一小节之后的 Galois 系数 \((K^s)^\times\) 将采用乘法记号。

\subsection{投射循环群的计算}
\label{b4:subsec:C02:02}

\begin{proposition}\label{b4:prop:C-procyclic-h1}
设 \(p\) 为素数。令 \(\sigma=1\) 表示加法群 \(\widehat{\vphantom{Z}\mathbb Z}\) 或 \(\mathbb Z_p\) 的拓扑生成元，并用 \(\sigma m\) 表示它对系数的作用，系数群均取离散拓扑。
\begin{enumerate}
\item 若 \(M\) 为有限离散连续 \(\widehat{\vphantom{Z}\mathbb Z}\)-模，则求值映射 \(f\mapsto f(\sigma)\) 诱导
\[
 H^1_{\mathrm{cts}}(\widehat{\vphantom{Z}\mathbb Z},M)
 \cong M/(\sigma-1)M.
\]
\item 若 \(M\) 为有限 \(p\)-主交换群，并有连续 \(\mathbb Z_p\) 作用，则同样有
\(H^1_{\mathrm{cts}}(\mathbb Z_p,M)\cong M/(\sigma-1)M\)。
\item 若 \(M\) 有限、\(|M|\) 与 \(p\) 互素，并有连续 \(\mathbb Z_p\) 作用，则 \(H^1_{\mathrm{cts}}(\mathbb Z_p,M)=0\)。
\end{enumerate}
\end{proposition}
\begin{proof}
余圈条件迫使 \(f(\sigma^r)=\sum_{i=0}^{r-1}\sigma^if(\sigma)\) 对 \(r>0\) 成立；逆元的值由 \(f(g^{-1})=-g^{-1}f(g)\) 决定。整数子群稠密，连续映射到 Hausdorff 离散空间由其在稠密子群上的值唯一决定，所以该求值映射是单射。

对第一项，令作用通过阶为 \(d\) 的循环商，取零化 \(M\) 的正整数 \(e\)。给定任意 \(a\in M\)，阶为 \(de\) 的循环群上的交叉同态可以用 \(f(\sigma)=a\) 定义，因为
\[
 \sum_{i=0}^{de-1}\sigma^ia
 =e\sum_{i=0}^{d-1}\sigma^ia=0.
\]
该等式正好保证绕循环一周后的关系成立；拉回到 \(\widehat{\vphantom{Z}\mathbb Z}\) 得到连续余圈。在第二项中，可取 \(d,e\) 都为 \(p\) 的幂，故相同构造经 \(\mathbb Z_p\) 的有限循环商完成。两种情形的主余圈在 \(\sigma\) 处取值恰为 \((\sigma-1)m\)，给出商空间公式。

第三项中，由\eqref{b4:eq:C-finite-quotient-cohomology}，一次余圈在某个有限循环 \(p\)-商 \(Q\) 上表示。令 \(S=\sum_{t\in Q}f(t)\)。对 \(g\in Q\)，按 \(gt\) 重编号得到 \(S=|Q|f(g)+gS\)。由于乘 \(|Q|\) 是 \(M\) 的自同构，可令 \(m=-S/|Q|\)，便得 \(f(g)=gm-m\)。所以所有类都为零。
\end{proof}

特别地，在平凡作用下，有限交换群 \(A\) 满足
\[
 H^1_{\mathrm{cts}}(\widehat{\vphantom{Z}\mathbb Z},A)\cong A,
 \qquad
 H^1_{\mathrm{cts}}(\mathbb Z_p,A)\cong A[p^\infty],
\]
其中 \(A[p^\infty]\) 是所有 \(p\)-幂阶元素组成的 \(p\)-主分量。第二式也可直接由连续同态的像必须为循环 \(p\)-群得到。

\begin{proposition}\label{b4:prop:C-procyclic-h2}
设 \(p\) 为素数，\(A\) 为有限交换群，赋予离散拓扑及所考察群的平凡作用。则
\[
 H^2_{\mathrm{cts}}(\widehat{\vphantom{Z}\mathbb Z},A)=0,
 \qquad H^2_{\mathrm{cts}}(\mathbb Z_p,A)=0.
\]
\end{proposition}
\begin{proof}
由\cref{b4:thm:cyclic-group-cohomology}，有限循环群 \(C_n\) 的平凡系数公式为 \(H^2(C_n,A)=A/nA\)。依据\cref{b4:thm:group-extensions-h2}，也可用中心扩张核验这一公式及其过渡映射：在 \(0\rightarrow A\rightarrow E\rightarrow C_n\rightarrow1\) 中，选生成元的提升 \(x\)，扩张由 \(x^n=a\in A\) 决定；把提升换成 \(bx\) 将 \(a\) 改成 \(a+nb\)。因为 \(A\) 居中而商群循环，\(A\) 与 \(x\) 生成整个 \(E\)，所以没有其他关系参数。

沿 \(C_{mn}\rightarrow C_n\) 拉回扩张后，新生成元提升的 \(mn\) 次幂为 \(ma\)，因此膨胀是
\[
 A/nA\longrightarrow A/(mn)A,\qquad [a]\longmapsto[ma].
\]
若 \(G=\widehat{\vphantom{Z}\mathbb Z}\)，取 \(m\) 为 \(A\) 的指数，每个类都在更后阶段变成零；\eqref{b4:eq:C-finite-quotient-cohomology}便给出所需消失。对 \(\mathbb Z_p\)，先分解有限群 \(A\) 的主分量。与 \(p\) 互素的部分已经满足 \(A/p^rA=0\)；\(p\)-主部分则可取 \(m\) 为足够大的 \(p\) 次幂以零化它。故同样消失。
\end{proof}

这里特意算了膨胀映射。各有限商的二次上同调通常不为零，例如 \(H^2(C_{p^r},\mathbb F_p)=\mathbb F_p\)，却在通向下一个 \(p\)-倍阶段时被送到零。非零有限阶段与零余极限完全相容。上述结论只涉及有限系数；不能据此把二次上同调对全部离散系数都宣布为零。

\subsection{Galois 上同调与局部域算例}
\label{b4:subsec:C02:03}

对域 \(K\) 及可分闭包 \(K^s\)，记 \(G_K=\operatorname{Gal}(K^s/K)\)。其开正规子群对应有限 Galois 子扩张。离散连续 \(G_K\)-模的上同调称为 \(K\) 的\term[Galoisshangtongdiao]{Galois 上同调}[Galois cohomology]，简记 \(H^n(K,M)\)。这里仍是连续理论。

\begin{theorem}[Hilbert 第 90 定理]\label{b4:thm:C-hilbert90}
设 \(K\) 为域，\(K^s\) 为可分闭包，\(G_K=\operatorname{Gal}(K^s/K)\) 取 Krull 拓扑。则
\[
 H^1_{\mathrm{cts}}\bigl(G_K,(K^s)^\times\bigr)=0.
\]
其中乘法群 \((K^s)^\times\) 取离散拓扑和自然 Galois 作用。
\end{theorem}
\begin{proof}
先设 \(L/K\) 为有限 Galois 扩张，群为 \(\Gamma\)，乘法一次余圈记作 \((c_\sigma)\)，满足 \(c_{\sigma\tau}=c_\sigma\sigma(c_\tau)\)。不同域自同构作为 \(L\)-值函数线性无关：若有项数最少的非零关系 \(\sum_i a_i\sigma_i=0\)，取 \(x\) 使 \(\sigma_1(x)\ne\sigma_r(x)\)，比较该关系在 \(xy\) 处的值与在 \(y\) 处的值乘 \(\sigma_r(x)\)，得到项数更少但非零的关系，矛盾；单项关系显然不可能。

因此可取 \(a\in L\)，使 \(b=\sum_{\sigma\in\Gamma}c_\sigma\sigma(a)\ne0\)。余圈公式及重编号给出
\[
 \tau(b)=\sum_{\sigma}\tau(c_\sigma)\tau\sigma(a)
 =c_\tau^{-1}\sum_\sigma c_{\tau\sigma}\tau\sigma(a)
 =c_\tau^{-1}b.
\]
令 \(d=b^{-1}\)，则 \(c_\tau=\tau(d)/d\)，所以余圈为主余圈。

一般连续余圈由\eqref{b4:eq:C-finite-quotient-cohomology}来自某个 \(G_K/U\) 与系数 \(\bigl((K^s)^\times\bigr)^U=L^\times\)，其中 \(L=(K^s)^U\) 为有限 Galois 扩张。有限情形使每个这样的类为零，故余极限也为零。
\end{proof}

这是 Hilbert 第 90 定理的乘法形式，连续版本及其有限层证明见 Weibel\cite[Hilbert's Theorem 90 6.11.16, pp. 213--214]{Weibel1994HomologicalAlgebra}。若正整数 \(n\) 在 \(K\) 中可逆（即特征为零，或正特征 \(p\) 不整除 \(n\)），\(K^s\) 中每个非零元素都有 \(n\) 次根，因而有离散连续模的 Kummer 正合列\footnote{库默尔（Ernst Eduard Kummer，1810-1893），德国数学家，研究代数数论，并提出理想数以处理分解问题。}
\[
 1\longrightarrow\mu_n\longrightarrow(K^s)^\times
 \xrightarrow{\;n\;}(K^s)^\times\longrightarrow1.
\]
由\cref{b4:prop:C-cochains-exact}的长正合列及\cref{b4:thm:C-hilbert90}，得到
\begin{equation}\label{b4:eq:C-kummer}
 H^1(K,\mu_n)\cong K^\times/(K^\times)^n.
\end{equation}
连接映射把 \(a\in K^\times\) 送到 \(g\mapsto g(b)/b\)，其中 \(b^n=a\)。换一个根只改变主余圈。系数 \(\mu_n\) 通常具有非平凡 Galois 作用，只有 \(K\) 已含全部 \(n\) 次单位根时，才可选根把它识别为平凡的 \(\mathbb Z/n\)-模。Kummer 列的这种系数区别也见 Weibel\cite[Chapter III, Sketch 6.10 and Example 6.10.1, pp. 241--242]{Weibel2013KBook}。

\begin{proposition}\label{b4:prop:C-kummer-local}
设 \(K=\mathbb F_q((t))\)，\(q\) 是素数 \(p\) 的幂，\(n\ge1\) 且 \(p\nmid n\)，\(\mu_n\) 为 \(K\) 的可分闭包中 \(n\) 次单位根组成的离散连续 Galois 模。则作为交换群，
\[
 H^1(K,\mu_n)\cong
 \mathbb Z/n\mathbb Z\oplus\mathbb Z/\gcd(n,q-1)\mathbb Z.
\]
\end{proposition}
\begin{proof}
每个非零 Laurent 级数唯一写成 \(t^r a u\)，其中 \(r\in\mathbb Z\)、\(a\in\mathbb F_q^\times\)、\(u\in1+t\mathbb F_q[\![t]\!]\)。在最后这个主单位群中，\(n\) 次幂映射为双射：要求
\((1+b_1t+b_2t^2+\cdots)^n=1+c_1t+c_2t^2+\cdots\)，第 \(r\) 次系数是 \(nb_r\) 加上只含 \(b_1,\ldots,b_{r-1}\) 的多项式。由于 \(n\) 在 \(\mathbb F_q\) 中可逆，每步唯一确定 \(b_r\)。因此幂商只留下 \(t^\mathbb Z\) 的 \(\mathbb Z/n\) 及有限循环群 \(\mathbb F_q^\times\) 的 \(n\) 次幂商，后者阶为 \(\gcd(n,q-1)\)。应用\eqref{b4:eq:C-kummer}即得结论。
\end{proof}

\(\mathbb F_q((t))\) 是非 Archimedes 局部域\footnote{阿基米德（\OriginalName{Ἀρχιμήδης}，约前287-前212），古希腊数学家、物理学家，研究几何求积与杠杆、浮力规律。}的一个基本模型：它对离散赋值完备，剩余域有限。上面的一次上同调使用了显式幂级数计算；二次上同调则与中心单代数的分裂有关。
\begin{theorem}[局部域的 Brauer 不变量]\label{b4:thm:C-local-brauer-invariant}
设 \(p\) 为素数，\(K\) 为 \(\mathbb Q_p\) 的有限扩张，\(K^s\) 为可分闭包，\(G_K=\operatorname{Gal}(K^s/K)\) 取 Krull 拓扑，\((K^s)^\times\) 取离散拓扑及自然 Galois 作用。则存在群同构
\[
 \operatorname{inv}_K:\operatorname{Br}(K)
 \cong H^2_{\mathrm{cts}}\bigl(G_K,(K^s)^\times\bigr)
 \xrightarrow{\sim}\mathbb Q/\mathbb Z.
\]
取归一化赋值 \(v_K\)。若 \(n\ge1\)，\(L/K\) 为次数 \(n\) 的非分歧扩张，\(\sigma\) 为算术 Frobenius 自同构\footnote{弗罗贝尼乌斯（Ferdinand Georg Frobenius，1849-1917），德国数学家，研究群表示、矩阵与代数数论。}，\(a\in K^\times\)，则循环代数 \((L/K,\sigma,a)\) 的不变量为 \(v_K(a)/n+\mathbb Z\)。
\end{theorem}
\begin{proofsketch}
第二卷第18.4节用交叉积证明有限 Galois 分裂扩张上的 Brauer 类\footnote{布劳尔（Richard Dagobert Brauer，1901-1977），德裔美国数学家，奠定模表示论，并研究中心单代数与类域论。}与乘法二次上同调对应。每个中心单代数都有有限可分分裂域，把它扩为 Galois 扩张，再取这些有限层的余极限，便得到 \(\operatorname{Br}(K)\cong H^2_{\mathrm{cts}}(G_K,(K^s)^\times)\)。

局部域论证的关键是：每个中心除代数 \(D/K\) 都被某个有限非分歧扩张分裂。设 \(\dim_KD=n^2\)，对非零 \(x\in D\) 定义
\[
 \nu(x)=\frac1n v_K(\operatorname{Nrd}(x)),\qquad \nu(0)=+\infty.
\]
在每个交换子域 \(K(x)\) 上，约化范数是该域范数的相应整数次幂，故 \(\nu\) 限制为完备赋值的唯一延拓。对 \(x,y\ne0\)，在交换域 \(K(x^{-1}y)\) 中使用 \(\nu(1+x^{-1}y)\ge\min(0,\nu(x^{-1}y))\)，再加 \(\nu(x)\)，得到三角不等式。约化范数的乘法性给出 \(\nu(xy)=\nu(x)+\nu(y)\)。由此得到 \(D\) 的赋值整环，且其拓扑是有限维 \(K\) 向量空间的通常拓扑。

若 \(n>1\)，\(D\) 必含一个非平凡非分歧交换子域。否则，每个 \(K(b)\subseteq D\) 的剩余域都只能等于 \(K\) 的剩余域：若不相等，剩余域扩张可分，Hensel 提升\footnote{亨泽尔（Kurt Hensel，1861-1941），德国数学家，引入 $p$ 进数并研究代数数论。}便在 \(K(b)\) 内给出所需非分歧子域。取 \(\pi\in D\) 的赋值为 \(\nu(D^\times)\) 的最小正值。任意整元素 \(b\) 可逐次写成
\[
 b=a_0+\pi a_1+\cdots+\pi^{N-1}a_{N-1}+\pi^Nb_N,
 \qquad a_i\in\mathcal O_K,\quad \nu(b_N)\ge0.
\]
余项趋于零，故 \(b\) 属于 \(K(\pi)\) 的闭包。有限维子空间 \(K(\pi)\) 已闭；再乘适当的 \(\pi\) 幂，把 \(D\) 中任意元素变成整元素，便会推出 \(D=K(\pi)\)，与中心为 \(K\)、\(n>1\) 矛盾。

现在对 \(n\) 归纳。取上述非分歧子域 \(K'/K\)，令 \(D'\) 为其中心化子。中心化子定理使 \(D'\) 是中心为 \(K'\) 的除代数，次数为 \(n/[K':K]<n\)。归纳给出 \(D'\) 中非分歧极大子域 \(L/K'\)；于是 \(L/K\) 仍非分歧，且 \([L:K]=n\)，是 \(D\) 的极大子域。极大子域分裂 \(D\)，证明关键步骤。这一直接证明见 Serre\cite[Chapter XII, Section 2, Proposition 1, Lemma 1 and Proposition 2, pp. 182--184]{Serre1979LocalFields}；有限交换域的赋值延拓及有限维范数拓扑见同书 Chapter II, Section 2。

设 \(L/K\) 为次数 \(n\) 的非分歧扩张。它的剩余域扩张为有限域扩张。剩余域的乘法范数满射，主单位群的逐层范数则由剩余域的迹给出；该迹满射。逐次修正单位，并用完备性取极限，得到
\(N_{L/K}(\mathcal O_L^\times)=\mathcal O_K^\times\)。另一方面，归一化赋值满足 \(v_K(N_{L/K}(x))=nv_L(x)\)。因此
\[
 K^\times/N_{L/K}(L^\times)\xrightarrow{\sim}\mathbb Z/n\mathbb Z,
 \qquad [a]\longmapsto v_K(a)\bmod n.
\]
循环群二次上同调公式和交叉积识别使左边等于 \(\operatorname{Br}(L/K)\)。

若把非分歧扩张的次数从 \(n\) 增到 \(mn\)，膨胀在赋值部分把 \([r]\) 送到 \([mr]\)。所以 \([r]\mapsto r/n+\mathbb Z\) 与有限层之间的映射相容。对全部正整数取余极限，得到 \(\mathbb Q/\mathbb Z\)；前述非分歧分裂保证没有遗漏任何 Brauer 类。循环代数的计算同时给出所列归一化。上同调与局部结构的关系可对照 Weibel\cite[Sections 6.11.17--6.11.19, pp. 214--215]{Weibel1994HomologicalAlgebra}；这项计算还需要除代数的局部结构，不能仅由一次 Kummer 公式推出。
\end{proofsketch}

% !TeX root = ../../Book4.tex
% 纲要标题：### C.3 层上同调、Čech 上同调与比较
% 纲要位置：工作记录/计划纲要.md，第 3918 行。
% 纲要细目（写作范围）：
% #### C.3.1 层与截面函子的右导出
% #### C.3.2 Čech 复形与覆盖计算
% #### C.3.3 零调覆盖与比较定理


\section{层上同调、Čech 上同调与比较}
\label{b4:sec:C03}

局部对象能够逐点存在，却未必能在整个空间上相容地选取。层上同调把这一障碍组织为截面函子的右导出；Čech 复形则直接记录一个覆盖上的拼接关系。它们的定义不同，比较需要假设。以下先给出交换群层的最小背景，再证明有限零调覆盖的比较定理。

\subsection{层与截面函子的右导出}
\label{b4:subsec:C03:01}

\begin{definition}\label{b4:def:C-sheaf}
设 \(X\) 为拓扑空间。给每个开集 \(U\) 一个交换群 \(\mathcal F(U)\)，给每个包含 \(V\subseteq U\) 一个限制同态 \(s\mapsto s|_V\)，并要求恒等限制与复合限制相容，便得到交换群预层。若对任意开覆盖 \(U=\bigcup_i U_i\)，在各 \(U_i\) 上给定的截面只要在全部交集上相同，就能唯一拼成 \(U\) 上的截面，则称这个预层为交换群\term[ceng]{层}[sheaf]。记全局截面为 \(\Gamma(X,\mathcal F)=\mathcal F(X)\)。
\end{definition}
\symidx{\(\Gamma(X,\mathcal F)\)：层的全局截面群}

层的态射是与所有限制同态相容的一族群同态。一个截面在点 \(x\) 附近的局部信息由其\term[jing]{茎}[stalk]
\[
 \mathcal F_x=\varinjlim_{x\in U}\mathcal F(U)
\]
记录；两截面给出同一元素，当且仅当它们在某个共同的更小邻域上相同。核逐开集计算，余核则需将逐开集的商作层化，即补入那些能局部表示且能相容拼接的截面。由此得到交换范畴 \(\mathbf{Sh}(X)\)。一列层正合，当且仅当它在每个茎上正合；特别地，层态射满射只要求截面局部能够提升，而不要求每个 \(\mathcal F(U)\rightarrow\mathcal G(U)\) 都满射。
\symidx{\(\mathcal F_x\)：层在一点的茎}
\symidx{\(\mathbf{Sh}(X)\)：拓扑空间上的交换群层范畴}

例如，在 \(\mathbb R\) 上把局部常值整数函数取在 \(0,1\) 两点的值，得到满层态射 \(\underline{\mathbb Z}_{\mathbb R}\to i_{0*}\mathbb Z\oplus i_{1*}\mathbb Z\)。但全局截面映射是 \(n\mapsto(n,n)\)，不能取得 \((1,0)\)：两个点附近可以分别取常数 \(1\) 与 \(0\)，而连通空间上的一个局部常值函数必须处处为同一常数。\cref{b4:prop:C-sheaf-global-lift-obstruction}核验这个反例。层的满射性保留局部提升的存在，却不能保证这些提升在全空间上相容；截面函子的右导出正是要记录这一障碍。

\begin{proposition}\label{b4:prop:C-sheaf-injectives}
对任意拓扑空间 \(X\)，交换群层范畴 \(\mathbf{Sh}(X)\) 有足够多内射对象。函子 \(\Gamma(X,-)\) 左正合。
\end{proposition}
\begin{proof}
截面函子保持核，因为一个截面映为零恰好表示它属于逐开集的核层，故左正合。要把一个层嵌入内射层，可以先用全部茎记录它的局部截面，再把每个茎嵌入内射交换群。对每个 \(x\in X\)，选单射 \(\iota_x:\mathcal F_x\hookrightarrow I_x\)，其中 \(I_x\) 为内射交换群。令 \(i_{x*}I_x\) 在包含 \(x\) 的开集上取 \(I_x\)，在其余开集上取零。它是层，且有自然同构
\[
 \operatorname{Hom}_{\mathbf{Sh}(X)}(\mathcal A,i_{x*}I_x)
 \cong\operatorname{Hom}_{\mathbb Z}(\mathcal A_x,I_x).
\]
给左侧态射在各邻域上的分量取余极限得到右侧；反过来，在包含 \(x\) 的开集上先取茎再施行给定同态，其余分量取零，构成逆映射。

茎函子正合，因为其滤过余极限保留局部核、像与商。因此 \(i_{x*}I_x\) 内射；这些对象的积也内射，因为到积的同态逐分量延拓。积层在开集上逐分量取积，所以自然态射 \(\mathcal F\rightarrow\prod_x i_{x*}I_x\) 在开集 \(U\) 上就是
\[
s\longmapsto\bigl(\iota_x(s_x)\bigr)_{x\in U},
\]
其中不属于 \(U\) 的点对应零群。若像为零，每个 \(s_x\) 均为零；这表示每个点附近都有一个开邻域使截面限制为零。层的唯一拼接条件便给出 \(s=0\)。因此该态射单射，得到所需内射嵌入。
\end{proof}

\begin{definition}\label{b4:def:C-sheaf-cohomology}
设 \(X\) 为拓扑空间，\(\mathcal F\) 为 \(X\) 上的交换群层。对 \(n\ge0\)，定义
\[
 H^n(X,\mathcal F)=R^n\Gamma(X,-)(\mathcal F).
\]
具体地，取内射消解 \(\mathcal F\rightarrow\mathcal I^\bullet\)，对截面复形 \(\Gamma(X,\mathcal I^\bullet)\) 取上同调。这称为 \(\mathcal F\) 的\term[cengshangtongdiao]{层上同调}[sheaf cohomology]。
\end{definition}
\symidx{\(H^n(X,\mathcal F)\)：层上同调}

上述内射构造与导出定义可对照 Weibel\cite[Example 2.3.12, p. 42; Application 2.5.4, p. 51]{Weibel1994HomologicalAlgebra}。层上同调按定义具有短正合列的长正合列，这与固定覆盖的 Čech 上同调有一个重要差别：后者的余链群由截面组成，层的满射未必使这些余链群满射。

\subsection{Čech 复形与覆盖计算}
\label{b4:subsec:C03:02}

\begin{definition}\label{b4:def:C-cech-complex}
设 \(\mathcal U=(U_1,\ldots,U_m)\) 为拓扑空间 \(X\) 的有限开覆盖，\(\mathcal F\) 为 \(X\) 上的交换群层。令 \(U_{i_0\cdots i_r}=U_{i_0}\cap\cdots\cap U_{i_r}\)，并对 \(r\ge0\) 定义
\[
 \check C^r(\mathcal U,\mathcal F)
 =\prod_{i_0<\cdots<i_r}\mathcal F(U_{i_0\cdots i_r}),
\]
交集为空时对应的截面群取零。微分为交错限制：
\[
 (\delta c)_{i_0\cdots i_{r+1}}
 =\sum_{j=0}^{r+1}(-1)^j
 c_{i_0\cdots\widehat{i_j}\cdots i_{r+1}}|_{U_{i_0\cdots i_{r+1}}}.
\]
这个复形称为该覆盖的\term[Cechfuxing]{Čech 复形}[Čech complex]，其上同调记为 \(\check H^r(\mathcal U,\mathcal F)\)。
\end{definition}
\symidx{\(\check C^r(\mathcal U,\mathcal F)\)：开覆盖上的 Čech 余链群}
\symidx{\(\check H^r(\mathcal U,\mathcal F)\)：给定开覆盖的 Čech 上同调}

连续删除两个指标时，两种删除次序的符号相反，故 \(\delta^2=0\)。零次余圈就是在两两交集上相同的一族局部截面，所以层公理给出 \(\check H^0(\mathcal U,\mathcal F)=\Gamma(X,\mathcal F)\)。这是零次的无条件识别，高次还须检查交集上的上同调。

\begin{example}\label{b4:exm:C-projective-line-cech}
设 \(k\) 为域，把 \(X=\mathbb P^1_k\) 写成两个仿射直线的并：\(U=\operatorname{Spec}k[t]\)、\(V=\operatorname{Spec}k[s]\)，交集上 \(s=t^{-1}\)。考虑在两开集上各为秩一自由层的 \(\mathcal L\)，选生成元 \(e_U,e_V\)，在交集上以 \(e_V=t^{-2}e_U\) 识别两个局部平凡化。这就是线丛层 \(\mathcal O_{\mathbb P^1}(-2)\)。两个局部截面分别写成 \(a(t)e_U\) 和 \(b(t^{-1})e_V\)。把后一项改用 \(e_U\) 表示，就成为 \(t^{-2}b(t^{-1})e_U\)，所以它们能拼接恰好要求 \(a(t)=t^{-2}b(t^{-1})\)。Čech 微分是交集上第二个限制减去第一个限制，因而复形为
\[
 k[t]\oplus k[t^{-1}]\longrightarrow k[t,t^{-1}],
 \qquad(a,b)\longmapsto t^{-2}b-a,
\]
其中目标已用 \(e_U\) 平凡化。\(k[t]\) 中的指数非负，\(t^{-2}k[t^{-1}]\) 中的指数至多为 \(-2\)，二者交为零。因此零次上同调为零；像恰由除 \(t^{-1}\) 外的全部 Laurent 单项式张成，故
\[
 \check H^1\bigl(\{U,V\},\mathcal L\bigr)\cong k\cdot[t^{-1}]\cong k,
 \qquad \check H^r\bigl(\{U,V\},\mathcal L\bigr)=0\quad(r\ge2).
\]
同类的双开集计算可对照 Hartshorne 在微分层上的 Example 4.0.3\cite[Chapter III, Example 4.0.3, pp. 219--220]{Hartshorne1977}。现在只完成了 Čech 计算；把它识别为层上同调还需要下一小节的比较依据。
\end{example}

对一般覆盖也可用同样的有序指标定义复形。覆盖细化通过限制产生上同调映射，不同细化指标选择给出同伦的余链映射；对全部覆盖取余极限所得 \(\check H^n(X,\mathcal F)\)，与一个固定覆盖的 \(\check H^n(\mathcal U,\mathcal F)\) 是两个记号。对任意拓扑空间 \(X\) 和交换群层 \(\mathcal F\)，一次的自然比较为同构，见 Hartshorne\cite[Chapter III, Exercise 4.4, p. 223]{Hartshorne1977}：
\[
\varinjlim_{\mathcal U}\check H^1(\mathcal U,\mathcal F)
\xrightarrow{\sim}H^1(X,\mathcal F).
\]
这条一般结论只涉及一次，不能据此把全部覆盖的 Čech 余极限在所有高次上都等同于导出层上同调。下面证明一个适用于各次的有限零调覆盖比较。
\symidx{\(\check H^n(X,\mathcal F)\)：关于开覆盖取余极限的 Čech 上同调群}

\subsection{零调覆盖与比较定理}
\label{b4:subsec:C03:03}

若一个层 \(\mathcal I\) 对任意开集包含 \(V\subseteq U\) 都使限制 \(\mathcal I(U)\rightarrow\mathcal I(V)\) 满射，则称它为\term[songchiceng]{松弛层}[flasque sheaf]。我们先说明内射层具有这一性质，以及它在 Čech 复形中的作用。

\ProseParagraphBreak
在双开集覆盖 \(X=U\cup V\) 中，给定交集上的截面 \(c\in\mathcal I(U\cap V)\)，松弛性使它延拓为 \(b\in\mathcal I(V)\)。于是 \(c=\delta(0,b)\)，所以一次 Čech 上同调为零。相容性中的误差能够延到一个覆盖开集，再通过修改该开集上的截面消去；一般有限覆盖的证明将重复这个操作。

\begin{lemma}\label{b4:lem:C-flasque-cech}
设 \(X\) 为拓扑空间。\(X\) 上的内射交换群层都是松弛层，限制到任何开子集后仍内射。若 \(\mathcal I\) 为 \(X\) 上的松弛交换群层，\(\mathcal U\) 为 \(X\) 的任意有限开覆盖，则增广 Čech 复形
\[
 0\longrightarrow\Gamma(X,\mathcal I)
 \longrightarrow\check C^0(\mathcal U,\mathcal I)
 \longrightarrow\check C^1(\mathcal U,\mathcal I)\longrightarrow\cdots
\]
正合。
此外，每个松弛交换群层 \(\mathcal I\) 都满足 \(H^q(X,\mathcal I)=0\)（\(q>0\)）。
\end{lemma}
\begin{proof}
对开包含 \(j:U\hookrightarrow X\)，延零函子 \(j_!\) 把 \(U\) 上的层延到 \(X\)。具体地，对开集 \(V\subseteq X\)，\((j_!\mathcal A)(V)\) 由满足如下条件的截面 \(s\in\mathcal A(V\cap U)\) 组成：每个 \(z\in V\setminus U\) 都有 \(V\) 内的开邻域 \(W\)，使 \(s|_{W\cap U}=0\)。等价地，\(s\) 的非零茎所在的支集在 \(V\) 中闭；闭性是相对于当前开集 \(V\) 而言。例如 \(U=(0,1)\subset\mathbb R\)，非零常整数截面不属于 \((j_!\underline{\mathbb Z}_U)(\mathbb R)\)，因为它在 \(0,1\) 附近都不为零；同一截面却属于 \((j_!\underline{\mathbb Z}_U)(U)\)，此时没有需要检查的外部点。于是茎在 \(U\) 内保持原值，在外部为零。这个构造满足
\(\operatorname{Hom}_X(j_!\mathcal A,\mathcal B)=\operatorname{Hom}_U(\mathcal A,\mathcal B|_U)\)，并且由茎的描述正合。所以其右伴随限制函子保持内射。

记 \(\underline{\mathbb Z}_U\) 为 \(U\) 上的局部常值整数层。给出从它到任一层的态射等价于指定常数 \(1\) 的像，故
\(\operatorname{Hom}_X(j_!\underline{\mathbb Z}_U,\mathcal I)=\mathcal I(U)\)。若 \(V\subseteq U\)，有自然单射 \(j_{V!}\underline{\mathbb Z}_V\rightarrow j_{U!}\underline{\mathbb Z}_U\)，在每个茎上为恒等或从零群出发。内射性使相应 Hom 限制满射，即 \(\mathcal I(U)\rightarrow\mathcal I(V)\) 满射。

下面对覆盖开集数 \(m\) 归纳证明正次 Čech 上同调为零。\(m=1\) 时没有正次余链。设 \(m>1\)，令
\[
Y=\bigcup_{i<m}U_i,
\qquad Y\cap U_m=\bigcup_{i<m}(U_i\cap U_m).
\]
\(\mathcal I|_Y\) 和 \(\mathcal I|_{Y\cap U_m}\) 都仍松弛，因为其中开集之间的限制就是 \(\mathcal I\) 原有的满射限制。取正次 \(r\) 的余圈 \(c\)。\((U_i)_{i<m}\) 是 \(Y\) 的覆盖，对 \(\mathcal I|_Y\) 使用归纳，再将所得余链放回原覆盖，便可减去一个余边，使所有不含指标 \(m\) 的分量为零。剩下的分量 \(c_{i_0\cdots i_{r-1}m}\) 在 \(Y\cap U_m\) 的覆盖 \((U_i\cap U_m)_{i<m}\) 上成为 \(r-1\) 次余圈。

若 \(r>1\)，对 \(\mathcal I|_{Y\cap U_m}\) 和这个有 \(m-1\) 个开集的覆盖再次用归纳，把这些分量写成余边，再将相应低一次余链放回含 \(m\) 的分量，即消去剩余部分。若 \(r=1\)，这些截面在交集上相同，故拼成 \(Y\cap U_m\) 上一个截面。松弛性使它延到 \(U_m\)；把此延拓作为零次余链的第 \(m\) 分量，其余分量为零，余边便是剩余部分。零次正合性和增广单射来自层公理，完成归纳。

最后证明松弛层对截面函子零调。给松弛层 \(\mathcal G\) 取内射嵌入，记商层为 \(\mathcal Q\)，得到 \(0\to\mathcal G\to\mathcal J\to\mathcal Q\to0\)。给定 \(\mathcal Q(U)\) 中的截面，它在一个开覆盖上能提升到 \(\mathcal J\)。若已在部分开集之并 \(W\) 上选好提升，要再加入开集 \(V\)，在 \(W\cap V\) 上，用 \(V\) 上的提升减去 \(W\) 上的提升，所得的差属于 \(\mathcal G(W\cap V)\)。将这个差延到 \(V\)，再从 \(V\) 上的提升中减去，两个提升就能拼接。把覆盖良序化并逐次作此修正，在极限阶段用层公理拼接，便得 \(\mathcal J(U)\to\mathcal Q(U)\) 满射。对开集包含 \(V\subseteq U\)，先把 \(\mathcal Q(V)\) 中的截面提升到 \(\mathcal J(V)\)，再利用内射层的松弛性延到 \(\mathcal J(U)\)，可知 \(\mathcal Q\) 也松弛。截面列的满射性和 \(\mathcal J\) 的内射性给出 \(H^1(X,\mathcal G)=0\)，且对 \(q\ge2\)，长正合列给出 \(H^q(X,\mathcal G)\cong H^{q-1}(X,\mathcal Q)\)。反复对松弛商层应用这一过程，全部正次上同调均为零。
\end{proof}

给定拓扑空间 \(X\) 上的交换群层 \(\mathcal F\)。若有限开覆盖 \(\mathcal U\) 中的每个非空有限交集 \(U\) 都满足 \(H^q(U,\mathcal F|_U)=0\)（\(q>0\)），就称 \(\mathcal U\) 为相对于 \(\mathcal F\) 的\term[lingtiao-fugai]{零调覆盖}[acyclic cover]。这个条件依赖所给的层，不是覆盖自身的绝对性质。

\begin{theorem}[零调覆盖的比较]\label{b4:thm:C-cech-comparison}
设 \(X\) 为拓扑空间，\(\mathcal F\) 为 \(X\) 上的交换群层，\(\mathcal U=(U_1,\ldots,U_m)\) 为 \(X\) 的有限开覆盖，记 \(U_{i_0\cdots i_r}=U_{i_0}\cap\cdots\cap U_{i_r}\)。若对每个非空有限交集 \(U_{i_0\cdots i_r}\) 及每个 \(q>0\)，都有
\[
 H^q(U_{i_0\cdots i_r},\mathcal F|_{U_{i_0\cdots i_r}})=0,
\]
则有自然同构
\[
 \check H^n(\mathcal U,\mathcal F)\cong H^n(X,\mathcal F)
 \qquad(n\ge0).
\]
\end{theorem}
\begin{proof}
取层内射消解 \(\mathcal F\rightarrow\mathcal I^\bullet\)，组成第一象限双复形
\[
 D^{r,q}=\prod_{i_0<\cdots<i_r}
 \mathcal I^q(U_{i_0\cdots i_r}).
\]
横向为 Čech 微分，纵向为内射消解微分，并在总微分中给后者乘 \((-1)^r\)。每条对角线有限，故可以使用\cref{b4:cor:double-complex-comparison}。

先沿横向取上同调。由\cref{b4:lem:C-flasque-cech}，每个内射层的增广 Čech 行正合，因此从 \(\Gamma(X,\mathcal I^\bullet)\) 到总复形的增广比较为拟同构，其上同调就是 \(H^n(X,\mathcal F)\)。这一步只用消解层的松弛性，并不要求原层 \(\mathcal F\) 在覆盖交集上零调。

再沿纵向考察：内射层限制到每个交集后仍内射，所以该列计算交集上的层上同调。现在才使用零调覆盖假设，使正次全部为零，零次恰为 \(\mathcal F(U_{i_0\cdots i_r})\)。因此从 \(\check C^\bullet(\mathcal U,\mathcal F)\) 到同一个总复形的增广比较也是拟同构。内射消解给出共同的总复形，而原层在交集上的高次消失才使原层的 Čech 复形能够直接替代它。两次比较给出所求同构；消解比较的同伦唯一性保证它与层态射相容，并不依赖消解选择。
\end{proof}

现在回到\cref{b4:exm:C-projective-line-cech}。所需仿射消失定理是：若 \(A\) 为 Noether 交换环，\(\mathcal F\) 为 \(X=\operatorname{Spec}A\) 上的\term[zhunningjuceng]{准凝聚层}[quasi-coherent sheaf]，即局部由模及其局部化给出的层，则 \(H^q(X,\mathcal F)=0\) 对 \(q>0\) 成立。下面给出其证明概要，具体模论细节可对照 Hartshorne\cite[Chapter III, Lemmas 3.2--3.3, Proposition 3.4 and Theorem 3.5, pp. 213--215]{Hartshorne1977}。
\begin{proofsketch}
这里使用仿射概形上准凝聚层与模的识别 \(\mathcal F=\widetilde M\)，以及 \(\widetilde M(D(f))=M_f\) 和 \(M\mapsto\widetilde M\) 的正合性。这些层论构造与识别见 Hartshorne\cite[Chapter II, Propositions 5.1--5.2 and 5.4, Corollary 5.5, pp. 110--113]{Hartshorne1977}；下文展开的是从内射模到松弛层的步骤。

写 \(\mathcal F=\widetilde M\)，并在 \(A\)-模范畴中取内射消解 \(M\to I^\bullet\)。局部化正合，使 \(\widetilde M\to\widetilde I^\bullet\) 为层消解；又有 \(\Gamma(X,\widetilde I^q)=I^q\)。若这些层对截面函子零调，取截面就重新得到原来的模消解，从而使正次上同调消失。因此要把模范畴中的内射性转成层的零调性；这里证明更强的结论：内射模 \(I\) 的伴随层 \(\widetilde I\) 松弛，再用\cref{b4:lem:C-flasque-cech}。其延拓过程是在基本开集 \(D(f)\) 上先处理截面，然后把剩余问题缩到 \(V(f)\)；下面两项模论事实分别保证这两步可以实行。

先固定内射模 \(I\)。对理想 \(\mathfrak a\)，令
\[
 J=\{u\in I:\mathfrak a^Nu=0\;\text{对某个}\;N\ge1\}.
\]
% 跨卷引用目标：b3:thm:artin-rees（定理13.2.1）、b2:thm:baer-criterion（定理7.2.3）。
它仍内射，这将使支于闭集 \(V(f)\) 的剩余截面继续满足归纳所需的内射模条件。事实上，对理想 \(\mathfrak b\) 到 \(J\) 的同态 \(\varphi\)，Noether 性使 \(\mathfrak b\) 有限生成，因而可取统一的 \(N\)，使 \(\varphi(\mathfrak a^N\mathfrak b)=0\)。第三卷定理13.2.1的 Artin--Rees 引理应用于有限生成模 \(A\)、子模 \(\mathfrak b\) 与理想 \(\mathfrak a\)，给出足够大的 \(r\)，使 \(\mathfrak b\cap\mathfrak a^r\subseteq\mathfrak a^N\mathfrak b\)。于是 \(\varphi\) 通过 \(\mathfrak b/(\mathfrak b\cap\mathfrak a^r)\) 分解。把它视为 \(A/\mathfrak a^r\) 的子模，利用 \(I\) 的内射性延拓到 \(A/\mathfrak a^r\)；延拓像被 \(\mathfrak a^r\) 零化，仍在 \(J\) 中。第二卷定理7.2.3的 Baer 判据证明所述内射性。

其次，\(I\to I_f\) 总满射，所以基本开集 \(D(f)\) 上的截面都能来自一个全局截面。Noether 性使理想链 \(\operatorname{ann}(f^r)\) 稳定，取稳定后的 \(r\)，对 \(u/f^n\) 定义
\[
 (f^{n+r})\longrightarrow I,\qquad af^{n+r}\longmapsto af^ru.
\]
稳定的零化子保证良定义。延拓到 \(A\) 后，若 \(1\) 的像为 \(v\)，则 \(f^{n+r}v=f^ru\)，故 \(v\) 在 \(I_f\) 中的像为 \(u/f^n\)。

固定 \(f\in A\)，把前面构造的 \(J\) 取为 \((f)\)-挠模。还须把支于 \(V(f)\) 的截面与 \(\widetilde J\) 联系起来。设 \(s\in\widetilde I(U)\) 在 \(U\cap D(f)\) 上限制为零。Noether 性使 \(U\) 有有限基本开覆盖 \(D(g_1),\ldots,D(g_l)\)。写
\[
s_i=s|_{D(g_i)}\in I_{g_i}.
\]
它在 \((I_{g_i})_f\) 中为零，由局部化的零判据，存在 \(N_i\) 使 \(f^{N_i}s_i=0\)。取这些指数的最大值 \(N\)，层的唯一性给出 \(f^Ns=0\) 于 \(U\)。若 \(s_i=u_i/g_i^{r_i}\)，再清除 \(g_i\) 的分母，得到某个 \(m_i\) 满足
\[
g_i^{m_i}f^Nu_i=0,
\qquad
s_i=\frac{g_i^{m_i}u_i}{g_i^{r_i+m_i}}\in J_{g_i}.
\]
这里分子被 \(f^N\) 零化，因而属于 \(J\)。这些局部截面在交集上相同，便拼成 \(\widetilde J(U)\) 的截面。反过来，\(J_f=0\)，所以 \(\widetilde J\) 的截面在 \(D(f)\) 上都为零。这证明了支于 \(V(f)\) 的截面子层正是 \(\widetilde J\)。

现在利用 \(X\) 的 Noether 性，对闭集 \(Y=\overline{\operatorname{Supp}\widetilde I}\) 作归纳：假定支集闭包包含于 \(Y\) 的真闭子集的所有内射模，其伴随层均已证明松弛。给定开集 \(U\) 上的截面，若 \(Y\cap U=\varnothing\)，它已为零。否则选取 \(D(f)\subseteq U\) 与 \(Y\) 相交。由刚证的满射，可用全局截面消去它在 \(D(f)\) 上的限制。余下截面支于 \(V(f)\)，由上一段正好属于 \(\widetilde J(U)\)。\(J\) 内射，且 \(\widetilde J\) 的支集闭包位于严格较小的 \(Y\cap V(f)\)；归纳使该截面也能延到全空间。因此每个开集上的截面都能全局延拓，\(\widetilde I\) 松弛。这里分别用了理想有限生成、零化子链稳定、Artin--Rees 以及闭集上的 Noether 归纳，这套论证的范围是 Noether 环。

因此 \(\widetilde I^\bullet\) 是截面函子的零调消解，由\cref{b4:thm:acyclic-resolution}可以计算 \(\widetilde M\) 的层上同调。取全局截面后重新得到原来的正合模消解 \(I^\bullet\)，其正次上同调为零，证明结论。
\end{proofsketch}

例中的 \(U,V,U\cap V\) 分别对应 Noether 环 \(k[t],k[t^{-1}],k[t,t^{-1}]\)，且 \(\mathcal L\) 在各处均为秩一自由层，满足该定理。因此\cref{b4:thm:C-cech-comparison}使已算得的结果成为
\[
 H^0\bigl(\mathbb P^1_k,\mathcal O(-2)\bigr)=0,
 \quad H^1\bigl(\mathbb P^1_k,\mathcal O(-2)\bigr)\cong k,
 \quad H^n\bigl(\mathbb P^1_k,\mathcal O(-2)\bigr)=0\ (n\ge2).
\]
若改用单开集覆盖 \(\{X\}\)，其 Čech 复形没有正次项，便会得到 \(\check H^1\bigl(\{X\},\mathcal L\bigr)=0\)。这说明固定覆盖必须满足比较条件；覆盖取得过粗，可能完全漏掉要计算的障碍。

\begin{proposition}\label{b4:prop:C-sheaf-global-lift-obstruction}
设 \(X=\mathbb R\) 取通常拓扑，\(\underline{\mathbb Z}_X\) 为局部常值整数函数层，\(i_0,i_1\) 分别为点 \(0,1\) 到 \(X\) 的包含。逐开集取这两点处的值，给出满层态射
\[
\underline{\mathbb Z}_X\longrightarrow i_{0*}\mathbb Z\oplus i_{1*}\mathbb Z,
\]
但其全局截面映射 \(\mathbb Z\to\mathbb Z\oplus\mathbb Z\) 不满射。因此全局截面函子一般不右正合。
\end{proposition}
\begin{proof}
取值与限制相容，因而得到层态射。在点 \(0\) 附近选避开 \(1\) 的小区间，目标层的茎为第一份 \(\mathbb Z\)，求值映射在该茎上为恒等。同理，点 \(1\) 的茎映射也是恒等。在其余点可选避开两点的邻域，目标层的茎为零，所以那里的茎映射也满射。层态射的满射性按茎检验，故所构造的态射满射。

\(X\) 连通，每个全局局部常值整数函数都是常数，所以源的全局截面群为 \(\mathbb Z\)。目标的全局截面群为 \(\mathbb Z\oplus\mathbb Z\)，诱导映射是 \(n\mapsto(n,n)\)，不能取得 \((1,0)\)。对该满层态射及其核取截面，短正合列的最后一个映射不满射，故截面函子不右正合。
\end{proof}

% !TeX root = ../../Book4.tex
% 纲要标题：### C.4 代数拓扑与平展上同调简介
% 纲要位置：工作记录/计划纲要.md，第 3629 行。
% #### C.4.3 纤维函子与 Galois 群重建
% 纲要细目（写作范围）：
% #### C.4.1 拓扑空间上的常值层
% #### C.4.2 平展上同调与 Galois 上同调
% ---


\section{代数拓扑与平展上同调简介}
\label{b4:sec:C04}

本节比较拓扑空间上的常值层上同调与域上的平展上同调。两者都由截面函子导出，但采用不同的覆盖和系数范畴。\secref{b4:sec:E04}从平展态射与覆盖定义相应的位点和层范畴，这里以域上的计算说明其同调意义。一般代数拓扑及 \(\ell\)-进上同调还需要进一步的比较与极限理论。

\subsection{拓扑空间上的常值层}
\label{b4:subsec:C04:01}

对拓扑空间 \(X\) 和交换群 \(A\)，开集 \(U\) 上所有局部常值函数 \(U\rightarrow A\) 组成一个层，记为 \(\underline A_X\)，称为\term[changzhiceng]{常值层}[constant sheaf]。若 \(U\) 不连通，它的截面未必是同一个常数；例如两个不相交区间上的截面为 \(A\oplus A\)。把每个非空开集都直接赋为 \(A\) 的常值预层通常还需要层化，不能忽略这一差别。

\ProseParagraphBreak
用两个连通开弧 \(U,V\) 覆盖圆周 \(S^1\)，使 \(U\cap V\) 有两个连通分支。常值层的 Čech 复形为
\[
 A\oplus A\longrightarrow A\oplus A,
 \qquad(a,b)\longmapsto(b-a,b-a).
\]
核为对角的 \(A\)，余核由 \((u,v)\mapsto v-u\) 识别为 \(A\)，故这个覆盖给出 \(\check H^0=A\)、\(\check H^1=A\)。当 \(A=\mathbb Z\) 时，这是 Hartshorne 的圆周计算\cite[Chapter III, Example 4.0.4, p. 220]{Hartshorne1977}。

\ProseParagraphBreak
把结果解释为通常的拓扑上同调，还需要比较定理。对有限单纯复形的几何实现及其开子集 \(X\)，常值层上同调与奇异上同调自然同构；当 \(X\) 本身是单纯复形的几何实现时，后者又与单纯上同调相同。相关比较背景见 Hartshorne\cite[Chapter III, Section 2, Historical Note, p. 211]{Hartshorne1977}。
\begin{proofsketch}
在每个开集 \(U\) 上取奇异余链复形 \(C^q(U;A)=\operatorname{Hom}_{\mathbb Z}(C_q(U),A)\)，再层化为 \(\mathcal S^q\)。常值函数给出增广 \(\underline A_X\to\mathcal S^0\)。\(X\) 有局部可缩的邻域基；在这些邻域中，正次奇异上同调为零，零次等于 \(A\)。茎是滤过余极限，而该余极限正合，故增广复形
\(0\to\underline A_X\to\mathcal S^0\to\mathcal S^1\to\cdots\) 在每个茎上正合，是层消解。

还要说明层化后的截面怎样由普通余链表示。任取开集 \(U\subseteq X\) 及 \(s\in\mathcal S^q(U)\)。由层化的定义，可取开覆盖 \((V_i)\) 及 \(c_i\in C^q(V_i;A)\)，使 \(c_i\) 在 \(V_i\) 上给出的层截面为 \(s|_{V_i}\)。\(U\) 可度量且仿紧，故可细化为局部有限开覆盖 \((W_i)\)，并使 \(\overline{W_i}\subseteq V_i\)，这里闭包在 \(U\) 中取。对每个奇异 \(q\)-单形 \(\sigma\)，若其像包含于某个 \(W_i\)，就选定一个这样的指标并令 \(c(\sigma)=c_i(\sigma)\)；若没有这样的指标，则令 \(c(\sigma)=0\)。奇异链群以这些单形为自由基，故这规定了普通余链 \(c\in C^q(U;A)\)。

固定 \(x\in U\)。局部有限性使某个邻域只与有限个 \(W_i\) 相交；对其中不满足 \(x\in\overline{W_i}\) 的指标，可再缩小邻域，避开它们的闭包。剩下的有限个指标均满足 \(x\in V_i\)。这些 \(c_i\) 在 \(x\) 处给出同一茎，因而可取共同邻域，使它们作为普通余链的限制全部相同。再将该邻域缩小为 \(N_x\)，使其包含于某个含 \(x\) 的 \(W_j\)。像位于 \(N_x\) 的单形至少可以选指标 \(j\)，而任何实际选中的指标都属于上述有限族，故取值总与 \(c_j\) 相同。因此 \(c|_{N_x}=c_j|_{N_x}\)，\(c\) 层化后在每个茎上恢复 \(s\)，从而恢复 \(s\) 本身。

这证明 \(C^q(U;A)\to\mathcal S^q(U)\) 对每个开集 \(U\) 都满射。若 \(V\subseteq U\)，先用 \(V\) 上的普通余链表示所给截面，再对像不完全位于 \(V\) 的奇异单形取零，将它延到 \(U\)。延拓的层截面限制回 \(V\) 就是原截面。因此 \(\mathcal S^q\) 松弛，可以用于计算层上同调。

一个普通余链层化后为零，当且仅当每点附近的限制为零，也就是它在某个开覆盖中的全部小单形上为零。记 \(C_*^{\mathcal U}(X)\) 为像包含在某个覆盖开集内的奇异单形生成的子复形。限制映射 \(C^q(X;A)\to\operatorname{Hom}_{\mathbb Z}(C_q^{\mathcal U}(X),A)\) 满射，因为可以在其余自由基单形上补零。结合刚证的层化满射及其核的描述，得到复形的自然同构
\[
 \Gamma(X,\mathcal S^\bullet)
 \cong\varinjlim_{\mathcal U}
 \operatorname{Hom}_{\mathbb Z}(C_*^{\mathcal U}(X),A),
\]
重心细分使每个单形经过有限次细分后变小；细分与恒等之间的棱柱同伦，逐维保持边界相容，给出 \(C_*^{\mathcal U}(X)\hookrightarrow C_*(X)\) 的链同伦等价。对任意交换群 \(A\) 取 Hom 后仍是余链同伦等价；这些比较与覆盖加细相容，滤过余极限正合，故上述余极限的上同调为奇异上同调。结合松弛消解便得到比较。当 \(X\) 是单纯复形的几何实现时，把单纯链嵌入奇异链并使用细分和逐维同伦，得到与单纯上同调的比较。

概要省略小链同伦在各维度的具体棱柱公式。局部可缩性用于茎上的正合性；每个开子集的仿紧性用于余链代表的局部有限拼接，因而也用于松弛性，不能对任意拓扑空间照搬。
\end{proofsketch}

圆周、开弧及其有限个区间分支满足上述条件，因而开弧上的正次常值层上同调为零。\cref{b4:thm:C-cech-comparison}遂给出 \(H^0(S^1,\underline A)=H^1(S^1,\underline A)=A\)，且 \(H^n(S^1,\underline A)=0\) 对 \(n>1\) 成立。一般空间仍须另行检查比较条件。

\subsection{平展上同调与 Galois 上同调}
\label{b4:subsec:C04:02}

代数几何中的 Zariski 开集\footnote{扎里斯基（Oscar Zariski，1899-1986），美国数学家，推动代数几何的代数化，并引入以其命名的拓扑。}有时不足以容纳所需局部构造。例如，在域 \(K\) 上添入一个可分多项式的根，通常要经过有限可分域扩张，不能只在单点空间 \(\operatorname{Spec}K\) 内缩小开集。平展覆盖将这种局部变换也纳入覆盖：在域的情形，基本对象是有限个有限可分扩张的直积，它们称为有限平展 \(K\)-代数。

\ProseParagraphBreak
一般的\term[pingzhanshangtongdiao]{平展上同调}[étale cohomology]以平展覆盖所定义的层范畴为基础，对全局截面取右导出。这里的对象与覆盖已超出通常拓扑空间的开子集，不应把同一个符号 \(H^n\) 理解为只改了下标的 Zariski 上同调。\secref{b4:sec:E04}给出小平展位点的定义，而\cref{b4:prop:field-finite-etale-basis}说明域上的层可在有限平展对象上计算。借助这一组基，交换群平展层对应离散连续 \(G_K\)-模 \(M\)，全局截面对应 \(M^{G_K}\)，所以
\begin{equation}\label{b4:eq:C-etale-field}
 H^n_{\mathrm{\acute et}}(\operatorname{Spec}K,M)
 \cong H^n_{\mathrm{cts}}(G_K,M).
\end{equation}
\begin{proofsketch}
固定可分闭包 \(K^s\)，记 \(G_K=\operatorname{Gal}(K^s/K)\)。对有限平展 \(K\)-代数 \(B\)，取有限集合 \(T_B=\operatorname{Hom}_{K\text{-}\mathrm{alg}}(B,K^s)\)，作用为 \(g\cdot t=g\circ t\)。每个传递有限连续 \(G_K\)-集为 \(G_K/H\)，其中 \(H\) 开；它对应有限可分域 \((K^s)^H\)。任意有限集分成有限个轨道，对应域的有限直积。具体地，对有限连续 \(G_K\)-集 \(T\)，令 \(B_T\) 为全部 \(G_K\)-等变函数 \(T\to K^s\) 组成的 \(K\)-代数，运算逐点进行。选一个有限 Galois 扩张同时包含上述各域及其全部嵌入，有限 Galois 对应给出典范同构
\[
 B\xrightarrow{\sim}B_{T_B},\qquad
 b\longmapsto\bigl(t\longmapsto t(b)\bigr).
\]
同样，\(t\mapsto(h\mapsto h(t))\) 给出等变同构 \(T\xrightarrow{\sim}T_{B_T}\)。
态射的两个方向也相互恢复：代数同态 \(u:B\to C\) 给出 \(T_C\to T_B\)、\(t\mapsto t\circ u\)；反过来，等变映射 \(a:T_C\to T_B\) 通过函数的拉回 \(h\mapsto h\circ a\) 给出 \(B_{T_B}\to B_{T_C}\)，由上述同构恢复唯一的 \(u\)。这些对应保持恒等与复合，故
\[
 \operatorname{Hom}_{K\text{-}\mathrm{alg}}(B,C)
 \cong\operatorname{Hom}_{G_K}(T_C,T_B).
\]
代数的这项反变等价与 \(\operatorname{Spec}\) 的反变方向相抵，得到有限平展 \(K\)-概形与有限连续 \(G_K\)-集的协变等价。以下把对应于概形 \(U\) 的集合记为 \(T(U)\)。平展覆盖对应集合上的联合满射。

给定离散连续 \(G_K\)-模 \(M\)，定义
\[
\begin{aligned}
 F_M(U)&=\operatorname{Hom}_{G_K}(T(U),M),\\
 F_M(f)(h)&=h\circ T(f)\quad(f:V\to U).
\end{aligned}
\]
联合满射的有限集覆盖上，等变函数唯一拼接，直接验证层公理。在 \(T(U)=G_K/H\) 时，取函数在单位陪集处的值，便将 \(F_M(U)\) 识别为 \(M^H\)。

反过来，给定交换群平展层 \(F\)，简记 \(F(B)=F(\operatorname{Spec}B)\)，沿 \(K^s/K\) 中的有限 Galois 扩张取余极限
\[
 M=\varinjlim_{N/K\ {\rm finite\ Galois}}F(N).
\]
\(g\in G_K\) 在 \(F(N)\) 上通过代数自同构 \(g|_N\) 的拉回作用；这些作用与过渡映射相容。每个元素来自某个 \(F(N)\)，被 \(G_N=\operatorname{Gal}(K^s/N)\) 固定，故作用连续。有限可分域扩张给出的态射是覆盖，层的分离性使所有过渡映射单射。对有限可分子扩张 \(L\subseteq K^s\)，记 \(G_L=\operatorname{Gal}(K^s/L)\)，选择有限 Galois \(N/K\) 包含 \(L\)。由
\(N\otimes_LN\cong\prod_{\sigma\in\operatorname{Gal}(N/L)}N\)，层的等化子条件给出 \(F(L)\cong F(N)^{\operatorname{Gal}(N/L)}\)。任一 \(m\in M^{G_L}\) 都可在某个同时包含 \(L\) 的有限 Galois 扩张 \(N/K\) 上取截面代表；该代表在 \(F(N)\) 中被 \(\operatorname{Gal}(N/L)\) 固定，因为 \(F(N)\to M\) 单射，而 \(G_L\to\operatorname{Gal}(N/L)\) 满射。因此
\[
 F(L)\xrightarrow{\sim}M^{G_L}.
\]

这项识别还要与全部态射相容。对 \(U=\operatorname{Spec}B\)，定义 \(\theta_U:F(U)\to F_M(U)\)：若 \(s\in F(U)\)、\(t:B\to K^s\)，取有限 Galois 扩张 \(N/K\) 包含 \(t(B)\)，令 \(\theta_U(s)(t)\) 为 \(s\) 沿 \(B\xrightarrow{t}N\) 拉回后在 \(M\) 中的类。增大 \(N\) 不改变这个类，且 \(\theta_U(s)(g\cdot t)=g\cdot\theta_U(s)(t)\)。在域的分量上，这就是刚证的固定元素识别；层将有限不交并送到截面群的直积，故每个 \(\theta_U\) 都是同构。对任意有限平展态射 \(f:V\to U\)，拉回的复合律给出
\[
 \theta_V\bigl(F(f)(s)\bigr)(t)
 =\theta_U(s)\bigl(T(f)(t)\bigr)
 \qquad(t\in T(V)),
\]
所以 \(\theta:F\xrightarrow{\sim}F_M\) 是层的自然同构。层态射在各 \(F(N)\) 上诱导相容映射，给出 \(G_K\)-模同态；模同态则通过函数的后复合给出层态射，上式保证这两个对应互逆。另一方面，从 \(M\) 出发再取几何茎得到 \(\varinjlim_N M^{G_N}=M\)，因为每个元素的开稳定子群包含一个开正规子群。由此两项构造在对象与态射上都互逆。

由\cref{b4:prop:field-finite-etale-basis}，限制到有限平展对象给出层范畴的等价，并与交换群运算相容。因此上述构造识别的是整个域上小平展位点的交换群层范畴。上述函子保持加法，给出交换范畴的等价，因而保持核、余核、短正合列及内射对象。全局截面对应 \(M^{G_K}\)，所以在两边选取相应内射消解，再施行截面或不变量函子，得到相同复形。这就证明\eqref{b4:eq:C-etale-field}。连续 Galois 系数及 Kummer 例子的描述亦见 Weibel\cite[Chapter III, Sketch 6.10 and Example 6.10.1, pp. 241--242]{Weibel2013KBook}。
\end{proofsketch}

一个完全可计算的例子是有限域 \(K=\mathbb F_q\)。其可分闭包是各有限扩张 \(\mathbb F_{q^r}\) 的并，Frobenius 映射 \(x\mapsto x^q\) 在每个有限层生成阶为 \(r\) 的 Galois 群，所以 \(G_K\cong\widehat{\vphantom{Z}\mathbb Z}\)。对平凡有限系数 \(A=\mathbb Z/n\mathbb Z\)，\cref{b4:prop:C-procyclic-h1}及\cref{b4:prop:C-procyclic-h2}给出
\[
 H^0_{\mathrm{\acute et}}(\operatorname{Spec}\mathbb F_q,A)=A,
 \quad H^1_{\mathrm{\acute et}}(\operatorname{Spec}\mathbb F_q,A)=A,
 \quad H^2_{\mathrm{\acute et}}(\operatorname{Spec}\mathbb F_q,A)=0.
\]
其 Zariski 底空间只有一个点，相应交换群层的截面函子却是恒等函子，正次上同调全为零。两者不同，原因在于覆盖和系数范畴不同，并非导出函子出了矛盾。

\ProseParagraphBreak
最后，\(\mu_n\) 与常值 \(\mathbb Z/n\) 的 Galois 作用也未必相同。\eqref{b4:eq:C-kummer}中必须保留 \(\mu_n\) 的自然作用；而 \(\ell\)-进系数 \(\mathbb Z_\ell=\varprojlim_r\mathbb Z/\ell^r\) 还涉及系数系统的逆极限。\chapref{b4:ch:13}已经说明逆极限未必正合，不能把有限系数的结论逐项写下后直接删除所有导出极限项。

\subsection{纤维函子与 Galois 群重建}
\label{b4:subsec:C04:03}

域的自同构可以置换每一个可分多项式的根，而且这些置换与根之间的代数关系相容。反过来，若知道全部有限平展对象的几何点及它们之间的映射，能否恢复这组自同构？先在有限连续群作用的模型中回答这个问题。

\begin{definition}\label{b4:def:galois-fiber-functor}
设 \(K\) 为域，固定可分闭包 \(K^s\)。从有限平展 \(K\)-概形到有限集合的函子
\[
 \omega_K:\mathrm{F\acute Et}(K)\longrightarrow\mathbf{FinSet},
 \qquad
 \operatorname{Spec}B\longmapsto
 \operatorname{Hom}_{K\text{-}\mathrm{alg}}(B,K^s)
\]
称为由这个可分闭包给出的\term[xianwei-hanzi]{纤维函子}[fiber functor]。对概形态射，它通过对应代数同态的前复合给出几何点之间的映射。
\end{definition}

固定几何点而后取纤维，才得到一个确定的函子。这个函子的自同构不是对每个有限集合任意选择置换：它们必须与全部有限平展态射同时相容。

\begin{theorem}\label{b4:thm:galois-fiber-reconstruction}
设 \(G\) 为投射有限群，\(\mathcal F_G\) 为有限连续左 \(G\)-集范畴，\(\omega:\mathcal F_G\to\mathbf{FinSet}\) 为忘却函子。则
\[
 G\xrightarrow{\sim}\operatorname{Aut}(\omega),
 \qquad
 g\longmapsto\bigl(t\mapsto g\cdot t\bigr)
\]
为群同构。给右侧赋予在各有限集合上逐点取值的拓扑，它还是拓扑群同构。自然变换可在 \(\mathcal F_G\) 的一个小骨架上计算。
\end{theorem}
\begin{proof}
等变映射与每个 \(g\) 的作用交换，所以所给分量自然，且相乘对应复合。设 \(\alpha\in\operatorname{Aut}(\omega)\)。对开正规子群 \(N\)，考虑正则左作用的有限集合 \(G/N\)。其右平移都是等变映射，故自然性使 \(\alpha_{G/N}\) 与全部右平移交换。若
\(g_N=\alpha_{G/N}(N)\)，就有
\[
 \alpha_{G/N}(xN)=g_NxN.
\]
对 \(N'\subseteq N\)，商映射 \(G/N'\to G/N\) 的自然性说明 \(g_{N'}\) 映到 \(g_N\)。于是 \((g_N)_N\) 是相容族，由 \(G\cong\varprojlim_NG/N\) 来自唯一的 \(g\in G\)。

任一有限连续 \(G\)-集 \(T\) 的作用都通过某个 \(G/N\)：有限多个点的开稳定子群之交为作用的开正规核。对 \(t\in T\)，映射
\[
 G/N\longrightarrow T,\qquad xN\longmapsto x\cdot t
\]
等变。将自然性用于它的单位陪集，得到 \(\alpha_T(t)=g\cdot t\)。这证明满射性；各正则商上的作用还立即给出单射性。

逐点拓扑的基本单位邻域是有限多个有限 \(G\)-集中指定点的稳定条件，其在 \(G\) 中的原像是有限个开稳定子群之交，故同构连续。反向看，要求在 \(G/N\) 的单位陪集上固定，恰好对应 \(N\)。这些 \(N\) 构成 \(G\) 的单位邻域基，故逆映射也连续。
\end{proof}

上一小节的范畴等价与纤维函子相容，因而给出
\[
 \operatorname{Gal}(K^s/K)\cong\operatorname{Aut}(\omega_K).
\]
有限域的情形尤其具体：\(\mathbb F_{q^r}\) 的几何点构成一个阶为 \(r\) 的循环正则集，各层之间的限制要求这些循环置换相容，恢复的群便是 \(\widehat{\vphantom{Z}\mathbb Z}\)。因此投射极限在这里不只是计算群的方式，它记录了一切有限层置换必须满足的相容关系。

\ProseParagraphBreak
与有限连续投射有限群集范畴等价的范畴，通常称为\term[Galois-fanchou]{Galois 范畴}[Galois category]。这里证明的是给定群作用模型及纤维函子以后的重建；从一套抽象公理推出这类模型，是更一般的 Galois 范畴定理。若更换纤维函子之间的识别，重建群的识别随之共轭，所以纤维的选择是重建资料的一部分。


\begin{chaptersummary}
连续群上同调与抽象群上同调使用相同的代数微分，却在余链空间和系数范畴上有本质区别。紧性使连续余链通过某个有限商，滤过余极限正合性再把它转成有限群的计算。过渡映射仍须实际求出：有限循环商的二次上同调可以非零，而它的类在更后阶段全部消失。

\ProseParagraphBreak
Galois 系数的作用携带算术信息。Hilbert 第 90 定理消去乘法一次上同调，Kummer 列于是把单位根系数的一次上同调化成幂商；局部域的赋值、剩余域单位和主单位分别贡献不同部分。把单位根模无条件换成常值模，或把离散系数直接换成 \(p\)-进拓扑系数，都会改变计算对象。

\ProseParagraphBreak
层上同调测量截面函子的非正合性，Čech 复形记录一个覆盖上的拼接。双复形比较说明，交集上的高次消失才使有限覆盖足以计算导出上同调。常值层、Zariski 层与平展层采用不同的局部结构；它们之间的联系依靠具体比较定理，不能仅凭名称相近便予以识别。

\ProseParagraphBreak
对投射有限群，全部有限连续作用还保存群本身：纤维函子的自然自同构在每个正则有限商上由一个群元素决定，相容性将它们拼成投射极限中的唯一元素。域的有限平展对象由此不仅提供上同调的系数范畴，也恢复绝对 Galois 群。
\end{chaptersummary}

\BookExercises

\begin{exercise}\label{b4:ex:C-open-kernel}
设 \(G\) 为投射有限群，\(H\) 为开子群。证明由左陪集作用得到的同态 \(G\rightarrow S_{G/H}\) 连续，并写出其核。说明为什么此核给出包含于 \(H\) 的开正规子群。
\end{exercise}
\begin{solution}
紧性使 \(G/H\) 有限。陪集 \(xH\) 的稳定子群为 \(xHx^{-1}\)，所以置换作用的核为 \(\bigcap_{x\in G}xHx^{-1}\)。只有有限多个不同共轭子群：正规化子包含 \(H\)，故共轭子群数不超过 \([G:H]\)。该有限交为开正规子群，并包含于取 \(x=1\) 的那项 \(H\)。有限离散目标中的同态核开即连续，得到结论。
\end{solution}

\begin{exercise}\label{b4:ex:C-no-common-kernel}
设 \(G\) 为投射有限群，\(M\) 为离散连续 \(G\)-模。证明任意有限多个元素 \(m_1,\ldots,m_s\) 都被某个开正规子群同时固定，并说明
\[
 M=\varinjlim_U M^U,
\]
其中 \(U\) 遍历开正规子群，按反向包含排列。若 \(M\) 作为交换群有限生成，证明整个作用通过一个有限连续商；此结论的逆命题是否成立？
\end{exercise}
\begin{solution}
每个 \(m_i\) 的稳定子群都开，有限交仍开；在该交中选取开正规子群 \(U\)，便可同时固定所有这些元素。若 \(V\subseteq U\)，则 \(M^U\subseteq M^V\)，而任意两个开正规子群都有共同的更小开正规子群，所以这些包含构成滤过系统。每个元素属于某个 \(M^U\)，其并就是 \(M\)，得到所列余极限。

若 \(m_1,\ldots,m_s\) 生成底交换群，固定这些生成元的 \(U\) 也固定它们的每个整数线性组合，因而固定全体 \(M\)。作用于是通过有限连续商 \(G/U\)。逆命题不成立：给 \(\mathbb Q\) 离散拓扑和任意 \(G\) 的平凡作用，该作用通过平凡商，但 \(\mathbb Q\) 作为交换群不有限生成。这也说明有限作用商的存在不要求系数群本身有限。
\end{solution}

\begin{exercise}\label{b4:ex:C-discontinuous-finite-quotient}
令 \(G=\prod_{r\ge1}\mathbb F_2\)，取积拓扑和加法群结构。构造一个到 \(\mathbb F_2\) 的抽象群满同态，其核指数为 \(2\) 却不开。可使用向量空间基的延拓。
\end{exercise}
\begin{solution}
有限支集子空间 \(D=\bigoplus_r\mathbb F_2\) 在积拓扑中稠密。全一向量 \(v=(1,1,\ldots)\) 不在 \(D\) 中，所以 \(v+D\) 是商向量空间 \(G/D\) 的非零元素。把它延拓成一组 \(\mathbb F_2\)-基，定义线性泛函在 \(v+D\) 上取 \(1\)，在其他基向量上取零。与商映射复合得到满同态 \(f:G\rightarrow\mathbb F_2\)，其核含 \(D\) 且指数为 \(2\)。若该核开，则也闭；闭集包含稠密的 \(D\) 就只能等于 \(G\)，与 \(f(v)=1\) 矛盾。因此这是不连续的抽象有限商。
\end{solution}

\begin{exercise}\label{b4:ex:C-invariants-not-exact}
令 \(G=\mathbb Z_p\) 经 \(C_p=\langle\sigma\rangle\) 作用于 \(M=\mathbb F_p[C_p]\)，\(J\) 为增广映射 \(M\to\mathbb F_p\) 的核。计算 \(J^G\)、\(H^1_{\mathrm{cts}}(G,J)\) 和 \(H^1_{\mathrm{cts}}(G,M)\)。证明连接映射 \(\mathbb F_p\to H^1_{\mathrm{cts}}(G,J)\) 为同构，而包含 \(J\hookrightarrow M\) 在一次上同调上诱导零映射。
\end{exercise}
\begin{solution}
令 \(z=\sigma-1\)。在特征 \(p\) 下，\(M\) 可写为 \(\mathbb F_p[z]/(z^p)\)，增广映射是取常数项，故 \(J=zM\)。作用的差 \(\sigma-1\) 就是乘 \(z\)，于是
\[
 J^G=\mathbb F_p z^{p-1},\qquad
 H^1_{\mathrm{cts}}(G,J)=J/zJ\cong\mathbb F_p,\qquad
 H^1_{\mathrm{cts}}(G,M)=M/zM\cong\mathbb F_p,
\]
后两式用到\cref{b4:prop:C-procyclic-h1}。用 \(1\in M\) 提升增广目标中的 \(1\)，连接余圈在 \(\sigma\) 处取值为 \(\sigma-1=z\)，其类生成 \(J/zJ\)，故连接映射是同构。包含映射把 \(z\) 的类送到 \(M/zM\) 中的零类，所以诱导的上同调映射为零。相比之下，增广映射在 \(M/zM\) 上诱导到 \(\mathbb F_p\) 的同构。这里底模的单射没有在一次上同调上保持为单射。
\end{solution}

\begin{exercise}\label{b4:ex:C-sign-action}
让 \(\widehat{\vphantom{Z}\mathbb Z}\) 的拓扑生成元在 \(M=\mathbb Z/8\mathbb Z\) 上按 \(-1\) 作用。求零次和一次连续上同调。这个作用能否由 \(\mathbb Z_3\) 的拓扑生成元给出？
\end{exercise}
\begin{solution}
作用通过 \(C_2\)，故连续。不变量是 \(2m=0\) 的元素 \(0,4\)，所以零次为 \(\mathbb Z/2\mathbb Z\)。\cref{b4:prop:C-procyclic-h1}给出一次为 \(M/2M\cong\mathbb Z/2\mathbb Z\)。若把同样的非平凡阶二自同构指定给 \(\mathbb Z_3\) 的拓扑生成元，就会要求一个连续满同态 \(\mathbb Z_3\rightarrow C_2\)，但 \(\mathbb Z_3\) 的有限商只有 \(3\)-幂阶，故不能这样定义连续作用。
\end{solution}

\begin{exercise}\label{b4:ex:C-unipotent-cochain-level}
设 \(p\) 为任意素数，\(G=\mathbb Z_p\)，\(M=\mathbb F_p e_1\oplus\mathbb F_p e_2\)，拓扑生成元 \(\sigma\) 满足 \(\sigma e_1=e_1\)、\(\sigma e_2=e_2+e_1\)。求一次连续上同调。对每个 \(a=u e_1+v e_2\)，写出由 \(f(\sigma)=a\) 决定的连续余圈，并在允许系数作用下降的商 \(C_{p^r}\)（\(r\ge1\)）中，求它所需的最小 \(r\)。比较 \(p=2\) 与奇素数的情形。
\end{exercise}
\begin{solution}
\((\sigma-1)M=\mathbb F_p e_1\)，所以\cref{b4:prop:C-procyclic-h1}给出 \(H^1_{\mathrm{cts}}(G,M)\cong\mathbb F_p\)，由 \(e_2\) 的类生成。余圈在非负整数幂处必须满足
\[
 f(\sigma^j)=\sum_{i=0}^{j-1}\sigma^i a
 =\left(ju+\binom j2v\right)e_1+jv e_2.
\]
这些公式要在 \(C_{p^r}\) 上成立，还须使 \(f(\sigma^{p^r})=0\)。作用本身通过 \(C_p\)，而
\[
 f(\sigma^p)=\binom p2v e_1.
\]
若 \(p\) 为奇素数，\(p\mid\binom p2\)，所以每个 \(a\) 都已给出 \(C_p\) 上的余圈，最小 \(r=1\)。若 \(p=2\)，该值为 \(v e_1\)：\(v=0\) 时仍取 \(r=1\)，\(v\ne0\) 时则须取 \(r=2\)，因为 \(f(\sigma^4)=0\)。将相应有限商上的余圈拉回到 \(G\)，便得连续余圈；整数幂的稠密性保证唯一性。上同调中的类由 \(v\) 决定，但余圈是否已在作用商上满足关系，还要另行检查。
\end{solution}

\begin{exercise}\label{b4:ex:C-nontorsion-h2}
给 \(\mathbb Z\)、\(\mathbb Q\) 和 \(\mathbb Q/\mathbb Z\) 平凡作用及离散拓扑。对 \(G=\widehat{\vphantom{Z}\mathbb Z}\) 或 \(G=\mathbb Z_p\)，利用短正合列 \(0\to\mathbb Z\to\mathbb Q\to\mathbb Q/\mathbb Z\to0\)，证明连接映射
\[
 \operatorname{Hom}_{\mathrm{cts}}(G,\mathbb Q/\mathbb Z)
 \xrightarrow{\sim}H^2_{\mathrm{cts}}(G,\mathbb Z).
\]
对有限循环商 \(C_n\)，把生成元送到 \(1/n+\mathbb Z\) 的特征提升为有理数值余链，写出它的连接二次余圈。
\end{exercise}
\begin{solution}
有限循环公式给出 \(H^1(C_n,\mathbb Q)=\mathbb Q[n]=0\) 和 \(H^2(C_n,\mathbb Q)=\mathbb Q/n\mathbb Q=0\)。\eqref{b4:eq:C-finite-quotient-cohomology}因此使 \(H^1_{\mathrm{cts}}(G,\mathbb Q)\) 与 \(H^2_{\mathrm{cts}}(G,\mathbb Q)\) 都为零。短正合列的长正合列便给出连接同构；平凡作用下的一次上同调正是连续同态群。

连续特征由拓扑生成元的像决定，该像须为有限阶元素。对 \(\widehat{\vphantom{Z}\mathbb Z}\)，每个 \(\mathbb Q/\mathbb Z\) 中元素都可作为像；对 \(\mathbb Z_p\)，像须为 \(p\)-幂阶元素，恰好组成 \(\mathbb Z[1/p]/\mathbb Z\)。这从连接映射重新得到\cref{b4:prop:C-integer-procyclic-h2}中的两群。

写 \(C_n=\langle\sigma\rangle\)，在 \(0\le i<n\) 时取提升 \(b(\sigma^i)=i/n\)。其余边为整数值二次余圈
\[
 (\mathrm db)(\sigma^i,\sigma^j)
 =b(\sigma^j)-b(\sigma^{i+j})+b(\sigma^i)
 =\left\lfloor\frac{i+j}{n}\right\rfloor.
\]
这里中间项的指数先模 \(n\) 化到 \(0,\ldots,n-1\)；越过该区间时留下的整数正是连接余圈。
\end{solution}

\begin{exercise}\label{b4:ex:C-hilbert90-complex}
对 \(\operatorname{Gal}(\mathbb C/\mathbb R)=\{1,\tau\}\)，其中 \(\tau\) 为复共轭，直接说明乘法一次余圈由满足 \(z\overline z=1\) 的 \(z\) 决定，并为每个这样的 \(z\) 构造 \(d\ne0\)，使 \(z=\overline d/d\)。
\end{exercise}
\begin{solution}
归一化给出 \(f(1)=1\)，余圈方程在 \((\tau,\tau)\) 处为 \(1=f(\tau)\overline{f(\tau)}\)，故只须指定这样的 \(z=f(\tau)\)。若 \(z\ne-1\)，取 \(d=1+\overline z\ne0\)，则
\(\overline d/d=(1+z)/(1+\overline z)=z\)，最后等式使用 \(z\overline z=1\)。若 \(z=-1\)，取 \(d=i\)，则 \(\overline d/d=-1\)。这给出本例所有余圈的显式主余圈表示。
\end{solution}

\begin{exercise}\label{b4:ex:C-local-kummer-examples}
求 \(K=\mathbb F_5(\!(t)\!)\) 的 \(H^1(K,\mu_2)\) 与 \(H^1(K,\mu_3)\)，并给出幂商代表。说明 \(\mu_3\) 是否为平凡 Galois 模。
\end{exercise}
\begin{solution}
\(\mathbb F_5^\times\) 为四阶循环群，平方商由 \(1,2\) 表示。因此平方商有代表 \(1,2,t,2t\)，群为 \((\mathbb Z/2\mathbb Z)^2\)。立方在四阶单位群上可逆，主单位的立方也可逆，故立方商只由赋值模 \(3\) 给出，代表为 \(1,t,t^2\)，群为 \(\mathbb Z/3\mathbb Z\)。\cref{b4:prop:C-kummer-local}或\eqref{b4:eq:C-kummer}将这些幂商识别为所求上同调。\(K\) 不含三次本原单位根：素于特征的有限阶单位在约化下保持阶，因为主单位中的相应有限阶元素只能为一；而 \(\mathbb F_5^\times\) 中没有三阶元素。因此 \(\mu_3\) 的 Galois 作用非平凡，不能把这里的系数直接写成常值 \(\mathbb Z/3\mathbb Z\)。
\end{solution}

\begin{exercise}\label{b4:ex:C-finite-field-kummer}
设 \(q\) 为素数 \(p\) 的幂，\(n\ge1\) 且 \(p\nmid n\)。分别用 Kummer 公式与投射循环群的一次上同调公式，计算 \(H^1(\mathbb F_q,\mu_n)\)，并检查两种答案相同。
\end{exercise}
\begin{solution}
Kummer 公式给出 \(\mathbb F_q^\times/(\mathbb F_q^\times)^n\)，这是阶为 \(\gcd(n,q-1)\) 的循环群。另一方面，绝对 Galois 群为 \(\widehat{\vphantom{Z}\mathbb Z}\)，其 Frobenius 生成元在 \(\mu_n\) 上按 \(\zeta\mapsto\zeta^q\) 作用。选本原根把底交换群识别为 \(\mathbb Z/n\mathbb Z\) 后，\(\sigma-1\) 对应乘 \(q-1\)。于是一次上同调为 \((\mathbb Z/n\mathbb Z)/(q-1)(\mathbb Z/n\mathbb Z)\)，同样是上述阶的循环群。这个比较同时保留了非平凡的系数作用。
\end{solution}

\begin{exercise}\label{b4:ex:C-inseparable-kummer-boundary}
令 \(K=\mathbb F_p(\!(t)\!)\)。证明 \(K^p=\mathbb F_p(\!(t^p)\!)\)，据此判断 \(1+t\) 是否有 \(p\) 次根，以及它是否能在 \(K\) 的可分闭包中取得 \(p\) 次根。再令 \(B=K[T]/(T^p-(1+t))\)，证明 \(B\) 是次数 \(p\) 的域扩张，且有限忠实平坦，但不是平展扩张。计算 \(B\otimes_K B\)。
\end{exercise}
\begin{solution}
在特征 \(p\) 中，Laurent 级数的 \(p\) 次幂满足
\[
 \left(\sum_{j\ge j_0}a_jt^j\right)^p
 =\sum_{j\ge j_0}a_jt^{pj},\qquad a_j\in\mathbb F_p.
\]
所以像中的非零项指数都被 \(p\) 整除；反过来，每个只含这类指数的级数都可把指数除以 \(p\) 取得根。这给出所述子域。\(1+t\) 虽然赋值为零，仍不是 \(p\) 次幂。其根在 \(K\) 上纯不可分，若还属于可分闭包，就必须属于 \(K\)，故也不可能在那里取得根。

在代数闭包中记唯一根为 \(\alpha\)。若 \(T^p-(1+t)\) 有次数 \(d\) 满足 \(0<d<p\) 的首一因子，该因子必为 \((T-\alpha)^d\)。其 \(T^{d-1}\) 系数为 \(-d\alpha\in K\)，而 \(d\) 在 \(K\) 中可逆，便推出 \(\alpha\in K\)，矛盾。因此该多项式不可约，\(B=K(\alpha)\) 的次数为 \(p\)。作为 \(K\)-模，\(B\) 自由且秩为 \(p>0\)，故忠实平坦；但扩张纯不可分，不是平展扩张。在 \(B\) 上令 \(U=T-\alpha\)，得到
\[
 B\otimes_K B\cong B[T]/((T-\alpha)^p)
 \cong B[U]/(U^p).
\]
这个基变换具有非零幂零元，也直接显示它不平展。添根在有限忠实平坦覆盖中可以完成，却不能用有限可分扩张完成，这与\cref{b4:prop:C-inseparable-kummer-boundary}的边界一致。
\end{solution}

\begin{exercise}\label{b4:ex:C-constant-presheaf}
设 \(X=U\sqcup V\)，其中 \(U,V\) 为非空连通开集，\(A\ne0\) 为交换群。说明把所有非空开集都赋为 \(A\)、非空之间限制均为恒等、在空集取零群且向空集的限制为零的预层为什么不是层，并求常值层在 \(X\) 上的截面。
\end{exercise}
\begin{solution}
在 \(U,V\) 上选两个不同元素。它们在空交集上的限制相同，但预层中的全局元素只能在两部分给出同一个值，不能拼出所选局部数据，违反存在性。常值层使用局部常值函数；连通性使每部分各有一个常数，而两部分之间没有相容条件，所以全局截面为 \(A\oplus A\)。
\end{solution}

\begin{exercise}\label{b4:ex:C-line-bundle-all-degrees}
把\cref{b4:exm:C-projective-line-cech}的粘合关系改为 \(e_V=t^d e_U\)，其中 \(d\in\mathbb Z\)，得到 \(\mathcal O_{\mathbb P^1}(d)\)。计算各次层上同调及其维数；允许使用本附录明确列出的仿射消失定理。
\end{exercise}
\begin{solution}
Čech 微分为 \((a,b)\mapsto t^d b-a\)。全局截面对应 \(k[t]\cap t^d k[t^{-1}]\)，其单项式指数满足 \(0\le j\le d\)，因此零次维数为 \(\max(d+1,0)\)。余核中未被像覆盖的指数满足 \(d<j<0\)，所以一次维数为 \(\max(-d-1,0)\)，可取这些 \(t^j\) 的类为基。覆盖只有两个开集，正次数大于一的 Čech 项为零。各交集仍仿射，限制仍为自由秩一层，因此仿射消失和\cref{b4:thm:C-cech-comparison}把这些结果全部识别为层上同调。
\end{solution}

\begin{exercise}\label{b4:ex:C-cover-refinement}
对 \(X=\mathbb P^1_k\)、\(\mathcal L=\mathcal O(-2)\)，比较单开集覆盖 \(\{X\}\) 与标准双开集覆盖 \(\{U,V\}\) 的一次 Čech 上同调。细化诱导的映射是否为同构？它是否说明对全部覆盖取余极限也一定错误？
\end{exercise}
\begin{solution}
单开集覆盖的一次 Čech 上同调为零；双开集覆盖的一次为 \(k\)。细化映射为 \(0\rightarrow k\)，故不是同构。这个例子只否定任取固定覆盖即可计算的说法，并不否定全部覆盖余极限；本例已有零调覆盖，所以层上同调确实为 \(k\)。覆盖的细化可能检出此前没有出现的类，不能将不同覆盖上的上同调群先验视为相同。
\end{solution}

\begin{exercise}\label{b4:ex:C-surjective-sheaf-map}
在 \(X=\mathbb R\) 上取 \(r\ge3\) 个不同点 \(x_1,\ldots,x_r\)，令 \(i_j:\{x_j\}\hookrightarrow X\)，并记求值满层态射
\[
 q:\underline{\mathbb Z}_X\longrightarrow\bigoplus_{j=1}^r i_{j*}\mathbb Z
\]
的核为 \(\mathcal K\)。计算 \(H^n(X,\mathcal K)\) 的各次群，确定连接映射在 \(\mathbb Z^r\) 上的核，并给出 \(H^1(X,\mathcal K)\) 的商群坐标。
\end{exercise}
\begin{solution}
各点的茎上只有相应求值分量，故 \(q\) 满射，与\cref{b4:prop:C-sheaf-global-lift-obstruction}的两点证明相同。每个 \(i_{j*}\mathbb Z\) 的限制映射都是恒等或到零群的映射，所以目标层松弛，其正次上同调为零。\(\mathbb R\) 同胚于开区间；\subsecref{b4:subsec:C04:01}的拓扑比较及可缩性也使常值整数层的正次上同调为零。

取核列的长正合列，得到
\[
 0\longrightarrow H^0(X,\mathcal K)\longrightarrow\mathbb Z
 \xrightarrow{\Delta}\mathbb Z^r
 \xrightarrow{\partial}H^1(X,\mathcal K)\longrightarrow0,
\]
且 \(H^n(X,\mathcal K)=0\) 对 \(n\ge2\) 成立。对角映射单射，所以 \(H^0(X,\mathcal K)=0\)。连接映射的核为 \(\Delta\mathbb Z\)，并诱导
\[
 H^1(X,\mathcal K)\cong\mathbb Z^r/\Delta\mathbb Z
 \cong\mathbb Z^{r-1},\qquad
 [(a_1,\ldots,a_r)]\longmapsto(a_2-a_1,\ldots,a_r-a_1).
\]
这些差值同时为零，恰好意味着所给各点取值能由一个全局常数实现。
\end{solution}

\begin{exercise}\label{b4:ex:C-flasque-not-injective}
证明松弛层不必是内射层。可取只有一个点的拓扑空间，并说明这个例子中所有层的正次层上同调仍为零。
\end{exercise}
\begin{solution}
单点空间上的交换群层范畴就是交换群范畴，每个层都松弛，因为只有到空集的零限制需要检查。但交换群 \(\mathbb Z\) 不内射：同态 \(2\mathbb Z\rightarrow\mathbb Z\)、\(2\mapsto1\) 不能延伸到 \(\mathbb Z\)，否则 \(1\) 的像须满足 \(2a=1\)。全局截面在此是恒等函子，因而正合，全部正次右导出均为零。零调性比内射性弱。
\end{solution}

\begin{exercise}\label{b4:ex:C-mayer-vietoris}
设 \(X=U\cup V\)，\(\mathcal F\) 为交换群层，不假定所给层 \(\mathcal F\) 在 \(U,V,U\cap V\) 上的正次层上同调消失。证明存在长正合列
\[
\begin{aligned}
 \cdots\longrightarrow H^n(X,\mathcal F)
 &\longrightarrow H^n(U,\mathcal F|_U)\oplus H^n(V,\mathcal F|_V)\\
 &\longrightarrow H^n(U\cap V,\mathcal F|_{U\cap V})\\
 &\longrightarrow H^{n+1}(X,\mathcal F)\longrightarrow\cdots.
\end{aligned}
\]
\end{exercise}
\begin{solution}
取内射消解 \(\mathcal F\rightarrow\mathcal I^\bullet\)。每项内射层松弛，所以\cref{b4:lem:C-flasque-cech}对双开集覆盖给出复形短正合列
\[
 0\longrightarrow\Gamma(X,\mathcal I^\bullet)
 \longrightarrow\Gamma(U,\mathcal I^\bullet)\oplus\Gamma(V,\mathcal I^\bullet)
 \longrightarrow\Gamma(U\cap V,\mathcal I^\bullet)\longrightarrow0.
\]
最后一个映射为第二个限制减去第一个限制。内射层在开集上的限制仍内射，所以三个截面复形分别计算题中的层上同调。对复形短正合列取长正合列即得结论；没有用到原层在各开集上的消失假设。
\end{solution}

\begin{exercise}\label{b4:ex:C-circle-finite-coefficients}
设 \(A=\mathbb Z/n\mathbb Z\)，用圆周的双开弧覆盖计算 \(H^1(S^1,\underline A)\)。写出一个代表生成元的 Čech 余圈，并解释改变局部零次截面怎样改变它。
\end{exercise}
\begin{solution}
交集有两个分支，余圈群为 \(A^2\)，余边是对角的 \(A\)。因此一次群为 \(A^2/\Delta A\cong A\)，由 \((u,v)\mapsto v-u\) 给出同构；\((0,1)\) 代表生成元。改变两个开弧上的零次截面 \(a,b\) 会把余圈改为 \((u+b-a,v+b-a)\)，不改变两个分支值之差。依据\subsecref{b4:subsec:C04:01}明确采用的拓扑比较，区间上的正次上同调为零，故该 Čech 计算确实给出层上同调。
\end{solution}

\begin{exercise}\label{b4:ex:C-two-etale-coefficients}
比较 \(\operatorname{Spec}\mathbb F_2\) 上常值 \(\mathbb Z/3\mathbb Z\) 系数与 \(\mu_3\) 系数的一次平展上同调。两系数的底交换群同构，为什么答案仍然不同？
\end{exercise}
\begin{solution}
绝对 Galois 群为 \(\widehat{\vphantom{Z}\mathbb Z}\)。常值系数的作用平凡，所以一次为 \(\mathbb Z/3\mathbb Z\)。在 \(\mu_3\) 上，Frobenius 为平方，即加法记号下乘 \(2\)，所以 \(\sigma-1\) 为乘 \(1\)，商群为零。也可由 Kummer 公式得到 \(\mathbb F_2^\times/(\mathbb F_2^\times)^3=1\)。不同之处是 Galois 作用；仅比较底群会丢失系数层的结构。
\end{solution}

\begin{exercise}\label{b4:ex:C-galois-cover-level}
设 \(K=\mathbb F_q\)，\(A=\mathbb Z/n\mathbb Z\) 为平凡系数，\(r\mid s\)。计算从覆盖 \(\mathbb F_{q^r}/K\) 加细到 \(\mathbb F_{q^s}/K\) 所诱导的一次 Čech 上同调映射。对 \(a\in A\)，求检测它所需的最小扩张次数；若从次数 \(r\) 的覆盖出发，求一个检测全部类的加细覆盖的最小次数。具体计算 \(n=12\)、\(r=8\) 的情形。
\end{exercise}
\begin{solution}
\cref{b4:prop:C-finite-galois-cover-h1}将两次覆盖的一次群分别识别为 \(A[r]\) 和 \(A[s]\)。膨胀把余圈沿 \(C_s\to C_r\) 拉回，生成元处的取值保持不变，所以映射是包含 \(A[r]\hookrightarrow A[s]\)。

若 \(a\) 的阶为 \(d\)，它属于 \(A[r]\) 当且仅当 \(d\mid r\)，因此检测该类的最小扩张次数为 \(d\)。检测全部类则要求 \(n\mid s\)，同时覆盖加细要求 \(r\mid s\)，故所求最小次数是 \(\operatorname{lcm}(r,n)\)。当 \(n=12\)、\(r=8\) 时，\(A[8]=\{0,3,6,9\}\) 只检出四个类；加细到次数 \(24\) 后，\(A[24]=A\)，十二个类全部出现。
\end{solution}

\begin{exercise}\label{b4:ex:C-fiber-finite-cyclic}
对有限循环群 \(C_m\)，从正则左作用及其右平移直接计算有限 \(C_m\)-集忘却函子的自然自同构群。说明为什么不能只检查一个平凡作用的集合。
\end{exercise}
\begin{solution}
自然自同构在正则集 \(C_m\) 上与每个右平移交换，所以由单位元的像 \(g\) 决定，且为左乘 \(g\)。对任意 \(C_m\)-集中的 \(t\)，等变映射 \(x\mapsto x\cdot t\) 又迫使该分量为 \(t\mapsto g\cdot t\)。每个 \(g\) 都给出这样的自然自同构，复合对应群乘法，因此答案为 \(C_m\)。平凡作用中每个 \(g\) 都作用为恒等，不能分辨群元素。
\end{solution}

\begin{exercise}\label{b4:ex:C-fiber-profinite-completion}
设 \(d\ge1\)，\(\Gamma=\mathbb Z^d\)，\(\omega\) 为有限左 \(\Gamma\)-集的忘却函子。证明 \(\operatorname{Aut}(\omega)\cong(\widehat{\vphantom{Z}\mathbb Z})^d\)，并说明自然映射 \(\Gamma\to\operatorname{Aut}(\omega)\) 单射但不满射。为证明不满射，显式构造一族相容剩余类：对每个 \(n\ge1\)，要求它在 \(n\) 的 \(2\) 次幂部分取零，在奇数部分取一。
\end{exercise}
\begin{solution}
每个有限指数子群 \(N\) 都包含某个 \(n\mathbb Z^d\)：可取 \(n=|\Gamma/N|\)，它零化整个有限商。因此子群 \(n\mathbb Z^d\) 在有限指数正规子群中共尾，\cref{b4:prop:C-abstract-fiber-completion}给出
\[
 \operatorname{Aut}(\omega)\cong\varprojlim_n(\mathbb Z/n\mathbb Z)^d
 \cong(\widehat{\vphantom{Z}\mathbb Z})^d.
\]
各坐标中的 \(\bigcap_n n\mathbb Z=0\)，故自然映射单射。

写 \(n=2^r m\)，其中 \(m\) 为奇数。中国剩余定理给出唯一的 \(a_n\in\mathbb Z/n\mathbb Z\)，满足
\[
 a_n\equiv0\pmod{2^r},\qquad a_n\equiv1\pmod m.
\]
当 \(n\mid n'\) 时，\(a_{n'}\) 模 \(n\) 的像仍满足这两条件，所以等于 \(a_n\)。若这族相容类来自整数 \(a\)，它在每个 \(2^r\) 模下为零，迫使 \(a=0\)；但它模 \(3\) 为 \(1\)，矛盾。于是 \((a_n)_n\) 不在 \(\mathbb Z\) 的像中，把它放在第一坐标、其余坐标取零，便得 \(d\) 维的不满射例子。全部有限作用可以区分原群元素，却仍可能增加不来自原群的自然自同构。
\end{solution}

\begin{exercise}\label{b4:ex:C-fiber-conjugation}
设 \(\omega,\omega'\) 为同一范畴的两个有限集合值函子，\(s:\omega\xrightarrow{\sim}\omega'\) 为自然同构。写出它诱导的自同构群同构，并说明更换 \(s\) 如何改变该同构。
\end{exercise}
\begin{solution}
所求为 \(\alpha\mapsto s\alpha s^{-1}\)，分量的自然性使结果仍为自然自同构。若另取 \(s'\)，令 \(\beta=s's^{-1}\in\operatorname{Aut}(\omega')\)，则 \(s'\alpha(s')^{-1}=\beta(s\alpha s^{-1})\beta^{-1}\)。因此两个群识别相差目标群中的内自同构。
\end{solution}
